%%%%%%%%%%%%%%%%%%%%%%%%%%%%%%%%%%%%%%%%%%%%%%%%%%%%%%%%%%%%
% 8.19 COMPUTATIONAL STATISTICS
% The Composition of Advanced Mathematics
%%%%%%%%%%%%%%%%%%%%%%%%%%%%%%%%%%%%%%%%%%%%%%%%%%%%%%%%%%%%

\section{Computational Statistics}
\label{sec:computational-statistics}

\prerequisite{
Probability, random variables, statistical estimation, hypothesis
testing, regression, mathematical statistics, Bayesian statistics,
statistical learning, numerical methods, optimization, linear algebra,
and basic programming concepts.
}

\learningobjectives{
By the end of this section, the reader should be able to:
\begin{itemize}
    \item explain why computation is fundamental to modern statistics;
    \item distinguish exact analytical methods from numerical methods;
    \item understand floating-point and numerical-precision issues;
    \item recognize numerical instability;
    \item understand numerical optimization in statistical estimation;
    \item recognize gradient, Hessian, and Newton-type methods;
    \item understand constrained optimization conceptually;
    \item understand Monte Carlo integration;
    \item understand Monte Carlo error and convergence rates;
    \item generate random variables using inverse transforms;
    \item understand acceptance--rejection sampling;
    \item understand importance sampling;
    \item recognize weight degeneracy;
    \item understand variance reduction conceptually;
    \item understand bootstrap resampling;
    \item distinguish parametric and nonparametric bootstrap methods;
    \item estimate bootstrap standard errors;
    \item understand percentile and basic bootstrap intervals;
    \item recognize bias-corrected bootstrap methods;
    \item understand jackknife resampling;
    \item understand permutation and randomization tests;
    \item distinguish exact and Monte Carlo permutation procedures;
    \item understand the plus-one correction for sampled permutations;
    \item understand Markov chain Monte Carlo;
    \item recognize Metropolis--Hastings;
    \item recognize Gibbs sampling;
    \item understand autocorrelation in MCMC;
    \item understand effective sample size;
    \item understand Monte Carlo standard errors;
    \item recognize importance sampling and sequential Monte Carlo;
    \item understand expectation--maximization;
    \item recognize latent-variable estimation;
    \item understand numerical differentiation;
    \item understand automatic differentiation conceptually;
    \item understand simulation studies;
    \item design computational experiments for estimator comparison;
    \item understand reproducibility and random seeds;
    \item understand parallel computation conceptually;
    \item recognize computational limitations in high-dimensional
          statistics;
    \item understand matrix decompositions used in statistical
          computation;
    \item understand iterative algorithms for regression and
          regularization;
    \item recognize coordinate descent;
    \item understand stochastic-gradient methods conceptually;
    \item understand computational diagnostics;
    \item connect computational statistics with numerical analysis,
          Bayesian statistics, statistical learning, and scientific
          computing.
\end{itemize}
}

%%%%%%%%%%%%%%%%%%%%%%%%%%%%%%%%%%%%%%%%%%%%%%%%%%%%%%%%%%%%
\subsection{What Is Computational Statistics?}
\label{subsec:what-is-computational-statistics}
%%%%%%%%%%%%%%%%%%%%%%%%%%%%%%%%%%%%%%%%%%%%%%%%%%%%%%%%%%%%

Computational statistics studies numerical algorithms for statistical
estimation, inference, simulation, prediction, and uncertainty
quantification.

Many statistical problems cannot be solved using closed-form formulas.

Instead, one must compute approximations to quantities such as

\[
\boxed{
\widehat{\theta}
=
\arg\max_\theta
L(\theta),
}
\]

\[
\boxed{
\mathbb{E}
[
g(\theta)
\mid x
],
}
\]

or

\[
\boxed{
P
(
T^*\geq T_{obs}
).
}
\]

%%%%%%%%%%%%%%%%%%%%%%%%%%%%%%%%%%%%%%%%%%%%%%%%%%%%%%%%%%%%
\subsection{Why Statistical Computation Is Necessary}
\label{subsec:why-statistical-computation}
%%%%%%%%%%%%%%%%%%%%%%%%%%%%%%%%%%%%%%%%%%%%%%%%%%%%%%%%%%%%

Computation becomes necessary when:

\[
\boxed{
\text{likelihood equations have no closed form},
}
\]

\[
\boxed{
\text{posterior integrals are intractable},
}
\]

\[
\boxed{
\text{sampling distributions are unknown},
}
\]

\[
\boxed{
\text{datasets are high dimensional},
}
\]

or

\[
\boxed{
\text{models require iterative estimation}.
}
\]

Modern statistics therefore combines theory with numerical algorithms.

%%%%%%%%%%%%%%%%%%%%%%%%%%%%%%%%%%%%%%%%%%%%%%%%%%%%%%%%%%%%
\subsection{Analytical Versus Computational Solutions}
\label{subsec:analytical-versus-computational}
%%%%%%%%%%%%%%%%%%%%%%%%%%%%%%%%%%%%%%%%%%%%%%%%%%%%%%%%%%%%

An analytical solution produces an exact symbolic formula.

For example,

\[
\boxed{
\widehat{\mu}
=
\overline{X}.
}
\]

A computational solution may instead produce a sequence

\[
\boxed{
\theta^{(0)},
\theta^{(1)},
\ldots
}
\]

that converges toward the desired estimate.

The output is then an approximation whose accuracy depends on the
algorithm and numerical precision.

%%%%%%%%%%%%%%%%%%%%%%%%%%%%%%%%%%%%%%%%%%%%%%%%%%%%%%%%%%%%
\subsection{Numerical Error}
\label{subsec:numerical-error-statistics}
%%%%%%%%%%%%%%%%%%%%%%%%%%%%%%%%%%%%%%%%%%%%%%%%%%%%%%%%%%%%

Computational error can arise from several sources:

\[
\boxed{
\text{roundoff error},
}
\]

\[
\boxed{
\text{truncation error},
}
\]

\[
\boxed{
\text{optimization error},
}
\]

\[
\boxed{
\text{Monte Carlo error},
}
\]

and

\[
\boxed{
\text{approximation error}.
}
\]

These are distinct from sampling error in the statistical model.

%%%%%%%%%%%%%%%%%%%%%%%%%%%%%%%%%%%%%%%%%%%%%%%%%%%%%%%%%%%%
\subsection{Sampling Error Versus Computational Error}
\label{subsec:sampling-versus-computational-error}
%%%%%%%%%%%%%%%%%%%%%%%%%%%%%%%%%%%%%%%%%%%%%%%%%%%%%%%%%%%%

Sampling error arises because only a finite random sample is observed.

Computational error arises because the statistical procedure is
implemented approximately.

Thus a reported estimate may contain both:

\[
\boxed{
\text{statistical uncertainty}
}
\]

and

\[
\boxed{
\text{numerical uncertainty}.
}
\]

A reliable analysis controls both.

%%%%%%%%%%%%%%%%%%%%%%%%%%%%%%%%%%%%%%%%%%%%%%%%%%%%%%%%%%%%
\subsection{Floating-Point Arithmetic}
\label{subsec:floating-point-statistics}
%%%%%%%%%%%%%%%%%%%%%%%%%%%%%%%%%%%%%%%%%%%%%%%%%%%%%%%%%%%%

Computers store most real numbers using finite binary representations.

Therefore many real numbers cannot be represented exactly.

For example,

\[
\boxed{
0.1
}
\]

does not generally have an exact finite binary representation.

Repeated calculations can accumulate small numerical errors.

%%%%%%%%%%%%%%%%%%%%%%%%%%%%%%%%%%%%%%%%%%%%%%%%%%%%%%%%%%%%
\subsection{Catastrophic Cancellation}
\label{subsec:catastrophic-cancellation}
%%%%%%%%%%%%%%%%%%%%%%%%%%%%%%%%%%%%%%%%%%%%%%%%%%%%%%%%%%%%

Catastrophic cancellation can occur when subtracting nearly equal
floating-point numbers.

For example,

\[
\boxed{
a-b
}
\]

may lose significant digits when \(a\) and \(b\) are very close.

Statistically equivalent formulas may therefore have very different
numerical stability.

%%%%%%%%%%%%%%%%%%%%%%%%%%%%%%%%%%%%%%%%%%%%%%%%%%%%%%%%%%%%
\subsection{Stable Variance Computation}
\label{subsec:stable-variance-computation}
%%%%%%%%%%%%%%%%%%%%%%%%%%%%%%%%%%%%%%%%%%%%%%%%%%%%%%%%%%%%

The identity

\[
\operatorname{Var}(X)
=
\mathbb{E}[X^2]
-
\mathbb{E}[X]^2
\]

is mathematically correct.

But computing sample variance as

\[
\frac1n\sum x_i^2
-
\overline{x}^2
\]

can suffer cancellation when the observations are large but tightly
clustered.

A centered calculation is generally more stable:

\[
\boxed{
\sum
(
x_i-\overline{x}
)^2.
}
\]

%%%%%%%%%%%%%%%%%%%%%%%%%%%%%%%%%%%%%%%%%%%%%%%%%%%%%%%%%%%%
\subsection{Logarithmic Computation}
\label{subsec:logarithmic-computation}
%%%%%%%%%%%%%%%%%%%%%%%%%%%%%%%%%%%%%%%%%%%%%%%%%%%%%%%%%%%%

Likelihoods often involve products of many small probabilities:

\[
\boxed{
L(\theta)
=
\prod_{i=1}^{n}
f_\theta(x_i).
}
\]

This product may underflow numerically.

Instead compute

\[
\boxed{
\ell(\theta)
=
\sum_{i=1}^{n}
\log
f_\theta(x_i).
}
\]

Log-likelihoods are both numerically more stable and easier to optimize.

%%%%%%%%%%%%%%%%%%%%%%%%%%%%%%%%%%%%%%%%%%%%%%%%%%%%%%%%%%%%
\subsection{Log-Sum-Exp Trick}
\label{subsec:log-sum-exp}
%%%%%%%%%%%%%%%%%%%%%%%%%%%%%%%%%%%%%%%%%%%%%%%%%%%%%%%%%%%%

Suppose one must compute

\[
\log
\left(
\sum_{i=1}^{n}
e^{a_i}
\right).
\]

Direct exponentiation may overflow.

Let

\[
m
=
\max_i a_i.
\]

Then

\[
\boxed{
\log
\left(
\sum_i e^{a_i}
\right)
=
m
+
\log
\left(
\sum_i
e^{a_i-m}
\right).
}
\]

This is known as the log-sum-exp stabilization.

%%%%%%%%%%%%%%%%%%%%%%%%%%%%%%%%%%%%%%%%%%%%%%%%%%%%%%%%%%%%
\subsection{Numerical Optimization in Statistics}
\label{subsec:numerical-optimization-statistics}
%%%%%%%%%%%%%%%%%%%%%%%%%%%%%%%%%%%%%%%%%%%%%%%%%%%%%%%%%%%%

Many estimators solve

\[
\boxed{
\widehat{\theta}
=
\arg\min_\theta
Q(\theta)
}
\]

or

\[
\boxed{
\widehat{\theta}
=
\arg\max_\theta
\ell(\theta).
}
\]

Examples include:

\[
\boxed{
\text{maximum likelihood},
}
\]

\[
\boxed{
\text{penalized regression},
}
\]

and

\[
\boxed{
\text{MAP estimation}.
}
\]

%%%%%%%%%%%%%%%%%%%%%%%%%%%%%%%%%%%%%%%%%%%%%%%%%%%%%%%%%%%%
\subsection{Gradient}
\label{subsec:statistical-gradient}
%%%%%%%%%%%%%%%%%%%%%%%%%%%%%%%%%%%%%%%%%%%%%%%%%%%%%%%%%%%%

For

\[
Q:
\mathbb{R}^p
\to
\mathbb{R},
\]

the gradient is

\[
\boxed{
\nabla Q(\boldsymbol{\theta})
=
\begin{pmatrix}
\partial Q/\partial\theta_1\\
\vdots\\
\partial Q/\partial\theta_p
\end{pmatrix}.
}
\]

At an interior differentiable optimum,

\[
\boxed{
\nabla Q(\widehat{\boldsymbol{\theta}})
=
\mathbf{0}.
}
\]

%%%%%%%%%%%%%%%%%%%%%%%%%%%%%%%%%%%%%%%%%%%%%%%%%%%%%%%%%%%%
\subsection{Hessian Matrix}
\label{subsec:statistical-hessian}
%%%%%%%%%%%%%%%%%%%%%%%%%%%%%%%%%%%%%%%%%%%%%%%%%%%%%%%%%%%%

The Hessian is

\[
\boxed{
H(\boldsymbol{\theta})
=
\nabla^2
Q(\boldsymbol{\theta}).
}
\]

It contains second derivatives.

At a local minimum, a positive definite Hessian indicates positive local
curvature.

For log-likelihood maximization, the negative Hessian is closely related
to observed Fisher information.

%%%%%%%%%%%%%%%%%%%%%%%%%%%%%%%%%%%%%%%%%%%%%%%%%%%%%%%%%%%%
\subsection{Gradient Descent}
\label{subsec:gradient-descent-statistics}
%%%%%%%%%%%%%%%%%%%%%%%%%%%%%%%%%%%%%%%%%%%%%%%%%%%%%%%%%%%%

Gradient descent iterates

\[
\boxed{
\boldsymbol{\theta}^{(k+1)}
=
\boldsymbol{\theta}^{(k)}
-
\eta_k
\nabla
Q
(
\boldsymbol{\theta}^{(k)}
),
}
\]

where

\[
\eta_k>0
\]

is the step size or learning rate.

The update moves in the direction of steepest local decrease.

%%%%%%%%%%%%%%%%%%%%%%%%%%%%%%%%%%%%%%%%%%%%%%%%%%%%%%%%%%%%
\subsection{Step Size}
\label{subsec:optimization-step-size}
%%%%%%%%%%%%%%%%%%%%%%%%%%%%%%%%%%%%%%%%%%%%%%%%%%%%%%%%%%%%

If the step size is too large, the algorithm may:

\[
\boxed{
\text{overshoot},
}
\]

\[
\boxed{
\text{oscillate},
}
\]

or

\[
\boxed{
\text{diverge}.
}
\]

If it is too small, convergence may be extremely slow.

Step-size selection is therefore an important computational issue.

%%%%%%%%%%%%%%%%%%%%%%%%%%%%%%%%%%%%%%%%%%%%%%%%%%%%%%%%%%%%
\subsection{Newton's Method}
\label{subsec:newtons-method-statistics}
%%%%%%%%%%%%%%%%%%%%%%%%%%%%%%%%%%%%%%%%%%%%%%%%%%%%%%%%%%%%

Newton's method uses both gradient and curvature:

\[
\boxed{
\boldsymbol{\theta}^{(k+1)}
=
\boldsymbol{\theta}^{(k)}
-
H
(
\boldsymbol{\theta}^{(k)}
)^{-1}
\nabla
Q
(
\boldsymbol{\theta}^{(k)}
).
}
\]

Near a well-behaved optimum, Newton's method can converge rapidly.

However, computing and inverting the Hessian may be expensive.

%%%%%%%%%%%%%%%%%%%%%%%%%%%%%%%%%%%%%%%%%%%%%%%%%%%%%%%%%%%%
\subsection{Newton--Raphson for Likelihood}
\label{subsec:newton-raphson-likelihood}
%%%%%%%%%%%%%%%%%%%%%%%%%%%%%%%%%%%%%%%%%%%%%%%%%%%%%%%%%%%%

For log-likelihood

\[
\ell(\theta),
\]

Newton--Raphson uses

\[
\boxed{
\theta^{(k+1)}
=
\theta^{(k)}
-
\frac{
\ell'(\theta^{(k)})
}{
\ell''(\theta^{(k)})
}.
}
\]

For vector parameters,

\[
\boxed{
\boldsymbol{\theta}^{(k+1)}
=
\boldsymbol{\theta}^{(k)}
-
[
\nabla^2\ell
(
\boldsymbol{\theta}^{(k)}
)
]^{-1}
\nabla\ell
(
\boldsymbol{\theta}^{(k)}
).
}
\]

%%%%%%%%%%%%%%%%%%%%%%%%%%%%%%%%%%%%%%%%%%%%%%%%%%%%%%%%%%%%
\subsection{Fisher Scoring}
\label{subsec:fisher-scoring}
%%%%%%%%%%%%%%%%%%%%%%%%%%%%%%%%%%%%%%%%%%%%%%%%%%%%%%%%%%%%

Fisher scoring replaces the observed negative Hessian by expected Fisher
information.

The update is

\[
\boxed{
\boldsymbol{\theta}^{(k+1)}
=
\boldsymbol{\theta}^{(k)}
+
I
(
\boldsymbol{\theta}^{(k)}
)^{-1}
U
(
\boldsymbol{\theta}^{(k)}
).
}
\]

This method is important in generalized linear models.

%%%%%%%%%%%%%%%%%%%%%%%%%%%%%%%%%%%%%%%%%%%%%%%%%%%%%%%%%%%%
\subsection{Iteratively Reweighted Least Squares}
\label{subsec:irls}
%%%%%%%%%%%%%%%%%%%%%%%%%%%%%%%%%%%%%%%%%%%%%%%%%%%%%%%%%%%%

Many generalized linear models can be fitted using iteratively reweighted
least squares, abbreviated IRLS.

At each iteration, the algorithm solves a weighted least-squares problem

\[
\boxed{
\min_{\boldsymbol{\beta}}
(
\mathbf{z}
-
X\boldsymbol{\beta}
)^T
W
(
\mathbf{z}
-
X\boldsymbol{\beta}
).
}
\]

The working response \(\mathbf{z}\) and weights \(W\) are updated
iteratively.

%%%%%%%%%%%%%%%%%%%%%%%%%%%%%%%%%%%%%%%%%%%%%%%%%%%%%%%%%%%%
\subsection{Convergence Criteria}
\label{subsec:optimization-convergence-criteria}
%%%%%%%%%%%%%%%%%%%%%%%%%%%%%%%%%%%%%%%%%%%%%%%%%%%%%%%%%%%%

An iterative algorithm may stop when:

\[
\boxed{
\|
\boldsymbol{\theta}^{(k+1)}
-
\boldsymbol{\theta}^{(k)}
\|
<
\varepsilon,
}
\]

or

\[
\boxed{
|
Q^{(k+1)}
-
Q^{(k)}
|
<
\varepsilon,
}
\]

or

\[
\boxed{
\|
\nabla Q
(
\boldsymbol{\theta}^{(k)}
)
\|
<
\varepsilon.
}
\]

Multiple criteria are often checked together.

%%%%%%%%%%%%%%%%%%%%%%%%%%%%%%%%%%%%%%%%%%%%%%%%%%%%%%%%%%%%
\subsection{Local Versus Global Optima}
\label{subsec:local-global-optima-statistics}
%%%%%%%%%%%%%%%%%%%%%%%%%%%%%%%%%%%%%%%%%%%%%%%%%%%%%%%%%%%%

A local optimum is best only in a neighborhood.

A global optimum is best over the full parameter space.

Nonconvex statistical models may contain several local optima.

Therefore initialization can materially affect the final estimate.

%%%%%%%%%%%%%%%%%%%%%%%%%%%%%%%%%%%%%%%%%%%%%%%%%%%%%%%%%%%%
\subsection{Multiple Starting Values}
\label{subsec:multiple-starting-values}
%%%%%%%%%%%%%%%%%%%%%%%%%%%%%%%%%%%%%%%%%%%%%%%%%%%%%%%%%%%%

For nonconvex problems, one can run an optimizer from several initial
points.

If many runs converge to the same optimum, confidence in the numerical
solution increases.

Different solutions may indicate:

\[
\boxed{
\text{multiple local optima},
}
\]

or

\[
\boxed{
\text{identifiability problems}.
}
\]

%%%%%%%%%%%%%%%%%%%%%%%%%%%%%%%%%%%%%%%%%%%%%%%%%%%%%%%%%%%%
\subsection{Constrained Optimization}
\label{subsec:constrained-optimization-statistics}
%%%%%%%%%%%%%%%%%%%%%%%%%%%%%%%%%%%%%%%%%%%%%%%%%%%%%%%%%%%%

Statistical parameters often satisfy constraints such as

\[
\boxed{
\sigma^2>0,
}
\]

\[
\boxed{
0<p<1,
}
\]

or

\[
\boxed{
\sum_jp_j=1.
}
\]

Optimization algorithms must respect these constraints.

One strategy is reparameterization.

%%%%%%%%%%%%%%%%%%%%%%%%%%%%%%%%%%%%%%%%%%%%%%%%%%%%%%%%%%%%
\subsection{Reparameterization}
\label{subsec:optimization-reparameterization}
%%%%%%%%%%%%%%%%%%%%%%%%%%%%%%%%%%%%%%%%%%%%%%%%%%%%%%%%%%%%

Instead of optimizing directly over

\[
\sigma>0,
\]

let

\[
\boxed{
\eta
=
\log\sigma.
}
\]

Then

\[
\boxed{
\sigma
=
e^\eta
}
\]

automatically satisfies positivity.

Similarly, probabilities can be represented using logistic or softmax
transformations.

%%%%%%%%%%%%%%%%%%%%%%%%%%%%%%%%%%%%%%%%%%%%%%%%%%%%%%%%%%%%
\subsection{Coordinate Descent}
\label{subsec:coordinate-descent}
%%%%%%%%%%%%%%%%%%%%%%%%%%%%%%%%%%%%%%%%%%%%%%%%%%%%%%%%%%%%

Coordinate descent updates one parameter at a time while holding the
others fixed.

For parameters

\[
\beta_1,\ldots,\beta_p,
\]

one cycles through

\[
\boxed{
\beta_1
\rightarrow
\beta_2
\rightarrow
\cdots
\rightarrow
\beta_p.
}
\]

Coordinate descent is especially useful for lasso and elastic-net
regression.

%%%%%%%%%%%%%%%%%%%%%%%%%%%%%%%%%%%%%%%%%%%%%%%%%%%%%%%%%%%%
\subsection{Soft Thresholding}
\label{subsec:soft-thresholding}
%%%%%%%%%%%%%%%%%%%%%%%%%%%%%%%%%%%%%%%%%%%%%%%%%%%%%%%%%%%%

The soft-thresholding operator is

\[
\boxed{
S(z,\lambda)
=
\operatorname{sign}(z)
\max
(
|z|-\lambda,
0
).
}
\]

It appears naturally in coordinate-descent updates for the lasso.

Values sufficiently close to zero are set exactly to zero.

%%%%%%%%%%%%%%%%%%%%%%%%%%%%%%%%%%%%%%%%%%%%%%%%%%%%%%%%%%%%
\subsection{Stochastic Gradient Descent}
\label{subsec:stochastic-gradient-descent}
%%%%%%%%%%%%%%%%%%%%%%%%%%%%%%%%%%%%%%%%%%%%%%%%%%%%%%%%%%%%

If

\[
Q(\theta)
=
\frac1n
\sum_{i=1}^{n}
Q_i(\theta),
\]

ordinary gradient descent computes the full gradient.

Stochastic gradient descent, abbreviated SGD, uses one observation or a
small batch:

\[
\boxed{
\theta^{(k+1)}
=
\theta^{(k)}
-
\eta_k
\nabla
Q_{i_k}
(
\theta^{(k)}
).
}
\]

This can dramatically reduce the cost per iteration for large datasets.

%%%%%%%%%%%%%%%%%%%%%%%%%%%%%%%%%%%%%%%%%%%%%%%%%%%%%%%%%%%%
\subsection{Mini-Batch Optimization}
\label{subsec:minibatch-optimization}
%%%%%%%%%%%%%%%%%%%%%%%%%%%%%%%%%%%%%%%%%%%%%%%%%%%%%%%%%%%%

A mini-batch method uses a subset

\[
B_k
\]

of observations at iteration \(k\):

\[
\boxed{
\theta^{(k+1)}
=
\theta^{(k)}
-
\eta_k
\frac1{|B_k|}
\sum_{i\in B_k}
\nabla Q_i
(
\theta^{(k)}
).
}
\]

This balances computational efficiency and gradient stability.

%%%%%%%%%%%%%%%%%%%%%%%%%%%%%%%%%%%%%%%%%%%%%%%%%%%%%%%%%%%%
\subsection{Numerical Differentiation}
\label{subsec:numerical-differentiation-statistics}
%%%%%%%%%%%%%%%%%%%%%%%%%%%%%%%%%%%%%%%%%%%%%%%%%%%%%%%%%%%%

If an analytic derivative is unavailable, a numerical approximation may
be used.

A forward difference is

\[
\boxed{
f'(x)
\approx
\frac{
f(x+h)-f(x)
}{
h
}.
}
\]

A centered difference is

\[
\boxed{
f'(x)
\approx
\frac{
f(x+h)-f(x-h)
}{
2h
}.
}
\]

The choice of \(h\) balances truncation and roundoff error.

%%%%%%%%%%%%%%%%%%%%%%%%%%%%%%%%%%%%%%%%%%%%%%%%%%%%%%%%%%%%
\subsection{Automatic Differentiation Preview}
\label{subsec:automatic-differentiation}
%%%%%%%%%%%%%%%%%%%%%%%%%%%%%%%%%%%%%%%%%%%%%%%%%%%%%%%%%%%%

Automatic differentiation computes derivatives by systematically applying
the chain rule through a sequence of elementary operations.

It is neither symbolic differentiation nor finite differencing.

It can produce derivatives accurate to machine precision and is widely
used in modern optimization and Bayesian computation.

%%%%%%%%%%%%%%%%%%%%%%%%%%%%%%%%%%%%%%%%%%%%%%%%%%%%%%%%%%%%
\subsection{Matrix Computation in Statistics}
\label{subsec:matrix-computation-statistics}
%%%%%%%%%%%%%%%%%%%%%%%%%%%%%%%%%%%%%%%%%%%%%%%%%%%%%%%%%%%%

Many statistical procedures require solving linear systems

\[
\boxed{
A\mathbf{x}
=
\mathbf{b}.
}
\]

Explicitly calculating

\[
A^{-1}
\]

is often less stable and less efficient than solving the system directly.

Thus numerical linear algebra is central to statistical computation.

%%%%%%%%%%%%%%%%%%%%%%%%%%%%%%%%%%%%%%%%%%%%%%%%%%%%%%%%%%%%
\subsection{QR Decomposition}
\label{subsec:qr-statistics}
%%%%%%%%%%%%%%%%%%%%%%%%%%%%%%%%%%%%%%%%%%%%%%%%%%%%%%%%%%%%

For a design matrix

\[
X,
\]

a QR decomposition writes

\[
\boxed{
X
=
QR,
}
\]

where \(Q\) has orthonormal columns and \(R\) is upper triangular.

Least squares can then be solved from

\[
\boxed{
R\widehat{\boldsymbol{\beta}}
=
Q^T\mathbf{y}.
}
\]

This is often more numerically stable than forming

\[
X^TX.
\]

%%%%%%%%%%%%%%%%%%%%%%%%%%%%%%%%%%%%%%%%%%%%%%%%%%%%%%%%%%%%
\subsection{Singular Value Decomposition}
\label{subsec:svd-statistics}
%%%%%%%%%%%%%%%%%%%%%%%%%%%%%%%%%%%%%%%%%%%%%%%%%%%%%%%%%%%%

The singular value decomposition is

\[
\boxed{
X
=
U
D
V^T.
}
\]

Here,

\[
U
\]

and

\[
V
\]

contain orthonormal singular vectors, and

\[
D
\]

contains nonnegative singular values.

SVD is fundamental for:

\[
\boxed{
\text{least squares},
}
\]

\[
\boxed{
\text{PCA},
}
\]

\[
\boxed{
\text{rank determination},
}
\]

and

\[
\boxed{
\text{regularization}.
}
\]

%%%%%%%%%%%%%%%%%%%%%%%%%%%%%%%%%%%%%%%%%%%%%%%%%%%%%%%%%%%%
\subsection{Condition Number}
\label{subsec:condition-number-statistics}
%%%%%%%%%%%%%%%%%%%%%%%%%%%%%%%%%%%%%%%%%%%%%%%%%%%%%%%%%%%%

For a nonsingular matrix \(A\), the condition number in the Euclidean norm
is

\[
\boxed{
\kappa(A)
=
\frac{
\sigma_{\max}(A)
}{
\sigma_{\min}(A)
}.
}
\]

A large condition number indicates that small perturbations can cause
large changes in the solution.

This is a major issue under multicollinearity.

%%%%%%%%%%%%%%%%%%%%%%%%%%%%%%%%%%%%%%%%%%%%%%%%%%%%%%%%%%%%
\subsection{Monte Carlo Methods}
\label{subsec:monte-carlo-methods-statistics}
%%%%%%%%%%%%%%%%%%%%%%%%%%%%%%%%%%%%%%%%%%%%%%%%%%%%%%%%%%%%

Monte Carlo methods approximate mathematical quantities through random
simulation.

Suppose

\[
X_1,\ldots,X_M
\overset{\text{i.i.d.}}{\sim}
p(x).
\]

To estimate

\[
\mu
=
\mathbb{E}
[
g(X)
],
\]

use

\[
\boxed{
\widehat{\mu}_M
=
\frac1M
\sum_{m=1}^{M}
g(X_m).
}
\]

%%%%%%%%%%%%%%%%%%%%%%%%%%%%%%%%%%%%%%%%%%%%%%%%%%%%%%%%%%%%
\subsection{Monte Carlo Consistency}
\label{subsec:monte-carlo-consistency}
%%%%%%%%%%%%%%%%%%%%%%%%%%%%%%%%%%%%%%%%%%%%%%%%%%%%%%%%%%%%

By the law of large numbers,

\[
\boxed{
\widehat{\mu}_M
\xrightarrow{P}
\mu.
}
\]

If

\[
\operatorname{Var}
[
g(X)
]
=
\sigma_g^2
<
\infty,
\]

then

\[
\boxed{
\operatorname{Var}
(
\widehat{\mu}_M
)
=
\frac{
\sigma_g^2
}{
M
}.
}
\]

%%%%%%%%%%%%%%%%%%%%%%%%%%%%%%%%%%%%%%%%%%%%%%%%%%%%%%%%%%%%
\subsection{Monte Carlo Standard Error}
\label{subsec:monte-carlo-se-independent}
%%%%%%%%%%%%%%%%%%%%%%%%%%%%%%%%%%%%%%%%%%%%%%%%%%%%%%%%%%%%

The Monte Carlo standard error is approximately

\[
\boxed{
MCSE
=
\frac{
s_g
}{
\sqrt{M}
},
}
\]

where \(s_g\) is the sample standard deviation of

\[
g(X_1),
\ldots,
g(X_M).
\]

Thus computational accuracy improves at rate

\[
\boxed{
M^{-1/2}.
}
\]

%%%%%%%%%%%%%%%%%%%%%%%%%%%%%%%%%%%%%%%%%%%%%%%%%%%%%%%%%%%%
\subsection{Monte Carlo Cost of Precision}
\label{subsec:monte-carlo-cost-precision}
%%%%%%%%%%%%%%%%%%%%%%%%%%%%%%%%%%%%%%%%%%%%%%%%%%%%%%%%%%%%

Because

\[
MCSE
\propto
M^{-1/2},
\]

reducing Monte Carlo error by a factor of \(2\) requires approximately

\[
\boxed{
4
}
\]

times as many simulations.

Reducing it by a factor of \(10\) requires approximately

\[
\boxed{
100
}
\]

times as many simulations.

%%%%%%%%%%%%%%%%%%%%%%%%%%%%%%%%%%%%%%%%%%%%%%%%%%%%%%%%%%%%
\subsection{Worked Example: Monte Carlo Integration}
\label{subsec:worked-monte-carlo-integration}
%%%%%%%%%%%%%%%%%%%%%%%%%%%%%%%%%%%%%%%%%%%%%%%%%%%%%%%%%%%%

\begin{workedexamplebox}{Estimating an Integral}

Suppose

\[
U
\sim
\operatorname{Uniform}(0,1).
\]

Then

\[
\int_0^1
x^2\,dx
=
\mathbb{E}
[
U^2
].
\]

Generate

\[
U_1,\ldots,U_M
\]

and calculate

\[
\widehat{I}
=
\frac1M
\sum_{m=1}^{M}
U_m^2.
\]

Since

\[
\int_0^1x^2\,dx
=
\frac13,
\]

the Monte Carlo estimate should approach

\[
\begin{answerbox}
\[
\boxed{
\frac13.
}
\]
\end{answerbox}

\end{workedexamplebox}

%%%%%%%%%%%%%%%%%%%%%%%%%%%%%%%%%%%%%%%%%%%%%%%%%%%%%%%%%%%%
\subsection{Pseudo-Random Numbers}
\label{subsec:pseudorandom-numbers}
%%%%%%%%%%%%%%%%%%%%%%%%%%%%%%%%%%%%%%%%%%%%%%%%%%%%%%%%%%%%

Computer-generated random numbers are generally produced by deterministic
algorithms designed to imitate randomness.

These are called pseudo-random number generators.

A generator begins from an internal state, often initialized by a seed.

%%%%%%%%%%%%%%%%%%%%%%%%%%%%%%%%%%%%%%%%%%%%%%%%%%%%%%%%%%%%
\subsection{Random Seeds}
\label{subsec:random-seeds}
%%%%%%%%%%%%%%%%%%%%%%%%%%%%%%%%%%%%%%%%%%%%%%%%%%%%%%%%%%%%

Setting a random seed can reproduce the same pseudo-random sequence.

This is useful for:

\[
\boxed{
\text{debugging},
}
\]

\[
\boxed{
\text{verification},
}
\]

and

\[
\boxed{
\text{reproducible research}.
}
\]

A fixed seed does not make the statistical procedure itself deterministic
in a conceptual sense; it fixes one reproducible simulation realization.

%%%%%%%%%%%%%%%%%%%%%%%%%%%%%%%%%%%%%%%%%%%%%%%%%%%%%%%%%%%%
\subsection{Uniform Random Variables}
\label{subsec:uniform-random-generation}
%%%%%%%%%%%%%%%%%%%%%%%%%%%%%%%%%%%%%%%%%%%%%%%%%%%%%%%%%%%%

Many random-generation algorithms begin with

\[
\boxed{
U
\sim
\operatorname{Uniform}(0,1).
}
\]

Other distributions are generated by transforming or accepting/rejecting
uniform random variables.

%%%%%%%%%%%%%%%%%%%%%%%%%%%%%%%%%%%%%%%%%%%%%%%%%%%%%%%%%%%%
\subsection{Inverse Transform Sampling}
\label{subsec:inverse-transform-sampling}
%%%%%%%%%%%%%%%%%%%%%%%%%%%%%%%%%%%%%%%%%%%%%%%%%%%%%%%%%%%%

Suppose \(F\) is a continuous strictly increasing CDF.

If

\[
U
\sim
\operatorname{Uniform}(0,1),
\]

then

\[
\boxed{
X
=
F^{-1}(U)
}
\]

has CDF \(F\).

This is inverse transform sampling.

%%%%%%%%%%%%%%%%%%%%%%%%%%%%%%%%%%%%%%%%%%%%%%%%%%%%%%%%%%%%
\subsection{Worked Example: Exponential Simulation}
\label{subsec:worked-exponential-simulation}
%%%%%%%%%%%%%%%%%%%%%%%%%%%%%%%%%%%%%%%%%%%%%%%%%%%%%%%%%%%%

\begin{workedexamplebox}{Inverse Transform Method}

For

\[
X
\sim
\operatorname{Exponential}(\lambda),
\]

the CDF is

\[
F(x)
=
1-e^{-\lambda x}.
\]

Set

\[
U
=
1-e^{-\lambda X}.
\]

Then

\[
e^{-\lambda X}
=
1-U.
\]

Thus

\[
X
=
-\frac1\lambda
\log(1-U).
\]

Since \(1-U\) is also uniform,

\begin{answerbox}
\[
\boxed{
X
=
-\frac1\lambda
\log U
}
\]

generates an exponential random variable.
\end{answerbox}

\end{workedexamplebox}

%%%%%%%%%%%%%%%%%%%%%%%%%%%%%%%%%%%%%%%%%%%%%%%%%%%%%%%%%%%%
\subsection{Acceptance--Rejection Sampling}
\label{subsec:acceptance-rejection}
%%%%%%%%%%%%%%%%%%%%%%%%%%%%%%%%%%%%%%%%%%%%%%%%%%%%%%%%%%%%

Suppose target density \(f\) is difficult to sample from.

Choose proposal density \(g\) and constant \(M\) such that

\[
\boxed{
f(x)
\leq
Mg(x)
}
\]

for all \(x\).

Generate

\[
Y\sim g
\]

and

\[
U\sim
\operatorname{Uniform}(0,1).
\]

Accept \(Y\) if

\[
\boxed{
U
\leq
\frac{
f(Y)
}{
Mg(Y)
}.
}
\]

%%%%%%%%%%%%%%%%%%%%%%%%%%%%%%%%%%%%%%%%%%%%%%%%%%%%%%%%%%%%
\subsection{Acceptance Rate}
\label{subsec:acceptance-rejection-rate}
%%%%%%%%%%%%%%%%%%%%%%%%%%%%%%%%%%%%%%%%%%%%%%%%%%%%%%%%%%%%

If \(f\) and \(g\) are normalized densities, the average acceptance
probability is

\[
\boxed{
\frac1M.
}
\]

Thus efficiency depends on choosing a proposal \(g\) that closely covers
the target.

A very large \(M\) produces many rejected proposals.

%%%%%%%%%%%%%%%%%%%%%%%%%%%%%%%%%%%%%%%%%%%%%%%%%%%%%%%%%%%%
\subsection{Importance Sampling}
\label{subsec:importance-sampling}
%%%%%%%%%%%%%%%%%%%%%%%%%%%%%%%%%%%%%%%%%%%%%%%%%%%%%%%%%%%%

Suppose

\[
\mu
=
\int
g(x)
p(x)\,dx
\]

but sampling from \(p\) is difficult.

Choose proposal density \(q\) with support covering the important regions
of \(p\).

Then

\[
\mu
=
\int
g(x)
\frac{
p(x)
}{
q(x)
}
q(x)\,dx.
\]

Therefore,

\[
\boxed{
\mu
=
\mathbb{E}_q
[
g(X)w(X)
],
}
\]

where

\[
\boxed{
w(x)
=
\frac{
p(x)
}{
q(x)
}.
}
\]

%%%%%%%%%%%%%%%%%%%%%%%%%%%%%%%%%%%%%%%%%%%%%%%%%%%%%%%%%%%%
\subsection{Importance Sampling Estimator}
\label{subsec:importance-sampling-estimator}
%%%%%%%%%%%%%%%%%%%%%%%%%%%%%%%%%%%%%%%%%%%%%%%%%%%%%%%%%%%%

If

\[
X_1,\ldots,X_M
\sim
q,
\]

then

\[
\boxed{
\widehat{\mu}
=
\frac1M
\sum_{m=1}^{M}
g(X_m)
w(X_m).
}
\]

This estimator works well only if \(q\) provides sufficient coverage of
important target regions.

%%%%%%%%%%%%%%%%%%%%%%%%%%%%%%%%%%%%%%%%%%%%%%%%%%%%%%%%%%%%
\subsection{Self-Normalized Importance Sampling}
\label{subsec:self-normalized-importance}
%%%%%%%%%%%%%%%%%%%%%%%%%%%%%%%%%%%%%%%%%%%%%%%%%%%%%%%%%%%%

If the target density is known only up to a constant,

\[
p(x)
\propto
\widetilde{p}(x),
\]

define unnormalized weights

\[
\widetilde{w}_m
=
\frac{
\widetilde{p}(X_m)
}{
q(X_m)
}.
\]

Normalize them:

\[
\boxed{
w_m
=
\frac{
\widetilde{w}_m
}{
\sum_{r=1}^{M}
\widetilde{w}_r
}.
}
\]

Then

\[
\boxed{
\widehat{\mu}
=
\sum_{m=1}^{M}
w_m
g(X_m).
}
\]

%%%%%%%%%%%%%%%%%%%%%%%%%%%%%%%%%%%%%%%%%%%%%%%%%%%%%%%%%%%%
\subsection{Importance Weight Degeneracy}
\label{subsec:importance-weight-degeneracy}
%%%%%%%%%%%%%%%%%%%%%%%%%%%%%%%%%%%%%%%%%%%%%%%%%%%%%%%%%%%%

If a few observations receive nearly all the importance weight, the
estimate becomes unstable.

This is called weight degeneracy.

It occurs when proposal \(q\) poorly matches the target.

Diagnostics should examine the distribution of weights.

%%%%%%%%%%%%%%%%%%%%%%%%%%%%%%%%%%%%%%%%%%%%%%%%%%%%%%%%%%%%
\subsection{Effective Sample Size for Importance Sampling}
\label{subsec:importance-sampling-ess}
%%%%%%%%%%%%%%%%%%%%%%%%%%%%%%%%%%%%%%%%%%%%%%%%%%%%%%%%%%%%

For normalized weights

\[
w_1,\ldots,w_M,
\]

a common heuristic effective sample size is

\[
\boxed{
ESS
=
\frac1{
\sum_{m=1}^{M}
w_m^2
}.
}
\]

If all weights are equal,

\[
ESS=M.
\]

If one weight dominates,

\[
ESS
\]

approaches \(1\).

%%%%%%%%%%%%%%%%%%%%%%%%%%%%%%%%%%%%%%%%%%%%%%%%%%%%%%%%%%%%
\subsection{Variance Reduction}
\label{subsec:monte-carlo-variance-reduction}
%%%%%%%%%%%%%%%%%%%%%%%%%%%%%%%%%%%%%%%%%%%%%%%%%%%%%%%%%%%%

Monte Carlo accuracy can be improved without simply increasing \(M\).

Variance-reduction methods include:

\[
\boxed{
\text{antithetic variables},
}
\]

\[
\boxed{
\text{control variates},
}
\]

\[
\boxed{
\text{importance sampling},
}
\]

and

\[
\boxed{
\text{stratified simulation}.
}
\]

%%%%%%%%%%%%%%%%%%%%%%%%%%%%%%%%%%%%%%%%%%%%%%%%%%%%%%%%%%%%
\subsection{Control Variates}
\label{subsec:control-variates}
%%%%%%%%%%%%%%%%%%%%%%%%%%%%%%%%%%%%%%%%%%%%%%%%%%%%%%%%%%%%

Suppose one wants

\[
\mu
=
\mathbb{E}[Y]
\]

and has another variable \(X\) with known mean

\[
\mu_X.
\]

Consider

\[
\boxed{
Y^*
=
Y
-
c
(
X-\mu_X
).
}
\]

Its expectation remains \(\mu\).

Choosing \(c\) appropriately can reduce variance when \(X\) and \(Y\) are
correlated.

%%%%%%%%%%%%%%%%%%%%%%%%%%%%%%%%%%%%%%%%%%%%%%%%%%%%%%%%%%%%
\subsection{Optimal Control-Variate Coefficient}
\label{subsec:optimal-control-variate}
%%%%%%%%%%%%%%%%%%%%%%%%%%%%%%%%%%%%%%%%%%%%%%%%%%%%%%%%%%%%

The variance-minimizing coefficient is

\[
\boxed{
c^*
=
\frac{
\operatorname{Cov}(Y,X)
}{
\operatorname{Var}(X)
}.
}
\]

Then

\[
\boxed{
\operatorname{Var}(Y^*)
=
\operatorname{Var}(Y)
(
1-\rho_{XY}^2
).
}
\]

Strong correlation can therefore yield substantial variance reduction.

%%%%%%%%%%%%%%%%%%%%%%%%%%%%%%%%%%%%%%%%%%%%%%%%%%%%%%%%%%%%
\subsection{Antithetic Variables}
\label{subsec:antithetic-variables}
%%%%%%%%%%%%%%%%%%%%%%%%%%%%%%%%%%%%%%%%%%%%%%%%%%%%%%%%%%%%

Suppose

\[
U
\sim
\operatorname{Uniform}(0,1).
\]

Then

\[
1-U
\]

has the same distribution.

A simulation may pair

\[
g(U)
\]

with

\[
g(1-U).
\]

If these outputs are negatively correlated, their average can have
smaller variance.

%%%%%%%%%%%%%%%%%%%%%%%%%%%%%%%%%%%%%%%%%%%%%%%%%%%%%%%%%%%%
\subsection{Bootstrap}
\label{subsec:bootstrap}
%%%%%%%%%%%%%%%%%%%%%%%%%%%%%%%%%%%%%%%%%%%%%%%%%%%%%%%%%%%%

The bootstrap approximates a sampling distribution by repeatedly
resampling from an estimated population distribution.

For a statistic

\[
\widehat{\theta}
=
T(X_1,\ldots,X_n),
\]

generate bootstrap datasets

\[
\boxed{
X_1^{*(b)},
\ldots,
X_n^{*(b)},
}
\]

and calculate

\[
\boxed{
\widehat{\theta}^{*(b)}.
}
\]

%%%%%%%%%%%%%%%%%%%%%%%%%%%%%%%%%%%%%%%%%%%%%%%%%%%%%%%%%%%%
\subsection{Nonparametric Bootstrap}
\label{subsec:nonparametric-bootstrap}
%%%%%%%%%%%%%%%%%%%%%%%%%%%%%%%%%%%%%%%%%%%%%%%%%%%%%%%%%%%%

The nonparametric bootstrap samples with replacement from the observed
dataset.

Each bootstrap observation is drawn from the empirical distribution

\[
\boxed{
F_n.
}
\]

Thus every original observation receives probability

\[
\boxed{
\frac1n
}
\]

on each bootstrap draw.

%%%%%%%%%%%%%%%%%%%%%%%%%%%%%%%%%%%%%%%%%%%%%%%%%%%%%%%%%%%%
\subsection{Bootstrap Algorithm}
\label{subsec:bootstrap-algorithm}
%%%%%%%%%%%%%%%%%%%%%%%%%%%%%%%%%%%%%%%%%%%%%%%%%%%%%%%%%%%%

A basic nonparametric bootstrap procedure is:

\begin{enumerate}
    \item observe data \(X_1,\ldots,X_n\);
    \item calculate \(\widehat{\theta}\);
    \item sample \(n\) observations with replacement from the data;
    \item calculate \(\widehat{\theta}^*\);
    \item repeat steps 3--4 for \(B\) bootstrap samples;
    \item use the empirical distribution of
          \(\widehat{\theta}^{*(1)},\ldots,\widehat{\theta}^{*(B)}\).
\end{enumerate}

%%%%%%%%%%%%%%%%%%%%%%%%%%%%%%%%%%%%%%%%%%%%%%%%%%%%%%%%%%%%
\subsection{Bootstrap Standard Error}
\label{subsec:bootstrap-standard-error}
%%%%%%%%%%%%%%%%%%%%%%%%%%%%%%%%%%%%%%%%%%%%%%%%%%%%%%%%%%%%

The bootstrap standard error is

\[
\boxed{
\widehat{SE}_{boot}
=
\sqrt{
\frac1{B-1}
\sum_{b=1}^{B}
\left(
\widehat{\theta}^{*(b)}
-
\overline{\widehat{\theta}^*}
\right)^2
}.
}
\]

Here,

\[
\boxed{
\overline{\widehat{\theta}^*}
=
\frac1B
\sum_{b=1}^{B}
\widehat{\theta}^{*(b)}.
}
\]

%%%%%%%%%%%%%%%%%%%%%%%%%%%%%%%%%%%%%%%%%%%%%%%%%%%%%%%%%%%%
\subsection{Worked Example: Bootstrap Concept}
\label{subsec:worked-bootstrap-concept}
%%%%%%%%%%%%%%%%%%%%%%%%%%%%%%%%%%%%%%%%%%%%%%%%%%%%%%%%%%%%

\begin{workedexamplebox}{Resampling a Small Dataset}

Suppose the observed sample is

\[
\{2,4,7\}.
\]

One bootstrap sample might be

\[
\{4,4,2\}.
\]

Another might be

\[
\{7,2,7\}.
\]

Each sample has size \(3\) and is drawn with replacement from the original
observations.

\begin{answerbox}
\[
\boxed{
\text{Bootstrap samples may contain repeated observations and omit
others.}
}
\]
\end{answerbox}

\end{workedexamplebox}

%%%%%%%%%%%%%%%%%%%%%%%%%%%%%%%%%%%%%%%%%%%%%%%%%%%%%%%%%%%%
\subsection{Bootstrap Bias Estimate}
\label{subsec:bootstrap-bias-estimate}
%%%%%%%%%%%%%%%%%%%%%%%%%%%%%%%%%%%%%%%%%%%%%%%%%%%%%%%%%%%%

A bootstrap estimate of estimator bias is

\[
\boxed{
\widehat{\operatorname{Bias}}_{boot}
=
\overline{\widehat{\theta}^*}
-
\widehat{\theta}.
}
\]

A simple bias-corrected point estimate is

\[
\boxed{
\widehat{\theta}_{BC}
=
\widehat{\theta}
-
\widehat{\operatorname{Bias}}_{boot}.
}
\]

Bias correction is not always beneficial because it can increase
variance.

%%%%%%%%%%%%%%%%%%%%%%%%%%%%%%%%%%%%%%%%%%%%%%%%%%%%%%%%%%%%
\subsection{Percentile Bootstrap Interval}
\label{subsec:percentile-bootstrap}
%%%%%%%%%%%%%%%%%%%%%%%%%%%%%%%%%%%%%%%%%%%%%%%%%%%%%%%%%%%%

Let

\[
\widehat{\theta}^{*}_{(\alpha/2)}
\]

and

\[
\widehat{\theta}^{*}_{(1-\alpha/2)}
\]

be empirical bootstrap quantiles.

The percentile interval is

\[
\boxed{
\left[
\widehat{\theta}^{*}_{(\alpha/2)},
\widehat{\theta}^{*}_{(1-\alpha/2)}
\right].
}
\]

This is simple but can have poor coverage in biased or highly skewed
settings.

%%%%%%%%%%%%%%%%%%%%%%%%%%%%%%%%%%%%%%%%%%%%%%%%%%%%%%%%%%%%
\subsection{Basic Bootstrap Interval}
\label{subsec:basic-bootstrap-interval}
%%%%%%%%%%%%%%%%%%%%%%%%%%%%%%%%%%%%%%%%%%%%%%%%%%%%%%%%%%%%

The basic bootstrap interval reflects the bootstrap distribution around
the observed estimate:

\[
\boxed{
\left[
2\widehat{\theta}
-
\widehat{\theta}^{*}_{(1-\alpha/2)},
\,
2\widehat{\theta}
-
\widehat{\theta}^{*}_{(\alpha/2)}
\right].
}
\]

It is sometimes called the reverse-percentile interval.

%%%%%%%%%%%%%%%%%%%%%%%%%%%%%%%%%%%%%%%%%%%%%%%%%%%%%%%%%%%%
\subsection{Bootstrap Normal Interval}
\label{subsec:bootstrap-normal-interval}
%%%%%%%%%%%%%%%%%%%%%%%%%%%%%%%%%%%%%%%%%%%%%%%%%%%%%%%%%%%%

Using bootstrap standard error,

\[
\boxed{
\widehat{\theta}
\pm
z_{1-\alpha/2}
\widehat{SE}_{boot}
}
\]

gives a normal-approximation bootstrap interval.

This assumes the estimator distribution is sufficiently symmetric and
approximately normal.

%%%%%%%%%%%%%%%%%%%%%%%%%%%%%%%%%%%%%%%%%%%%%%%%%%%%%%%%%%%%
\subsection{BCa Bootstrap Preview}
\label{subsec:bca-bootstrap-preview}
%%%%%%%%%%%%%%%%%%%%%%%%%%%%%%%%%%%%%%%%%%%%%%%%%%%%%%%%%%%%

The bias-corrected and accelerated, or BCa, interval adjusts bootstrap
quantiles for:

\[
\boxed{
\text{median bias}
}
\]

and

\[
\boxed{
\text{skewness-related acceleration}.
}
\]

It often improves coverage but requires additional computation.

%%%%%%%%%%%%%%%%%%%%%%%%%%%%%%%%%%%%%%%%%%%%%%%%%%%%%%%%%%%%
\subsection{Parametric Bootstrap}
\label{subsec:parametric-bootstrap}
%%%%%%%%%%%%%%%%%%%%%%%%%%%%%%%%%%%%%%%%%%%%%%%%%%%%%%%%%%%%

The parametric bootstrap assumes a fitted parametric model

\[
F_{\widehat{\theta}}.
\]

Bootstrap samples are generated from

\[
\boxed{
F_{\widehat{\theta}}
}
\]

rather than directly from the empirical distribution.

This can be more efficient when the parametric model is correct.

%%%%%%%%%%%%%%%%%%%%%%%%%%%%%%%%%%%%%%%%%%%%%%%%%%%%%%%%%%%%
\subsection{Parametric Versus Nonparametric Bootstrap}
\label{subsec:parametric-versus-nonparametric-bootstrap}
%%%%%%%%%%%%%%%%%%%%%%%%%%%%%%%%%%%%%%%%%%%%%%%%%%%%%%%%%%%%

The nonparametric bootstrap relies less on parametric distributional
assumptions.

The parametric bootstrap can model features not well represented by a
small empirical sample, but its validity depends strongly on model
specification.

%%%%%%%%%%%%%%%%%%%%%%%%%%%%%%%%%%%%%%%%%%%%%%%%%%%%%%%%%%%%
\subsection{Bootstrap for Regression}
\label{subsec:bootstrap-regression}
%%%%%%%%%%%%%%%%%%%%%%%%%%%%%%%%%%%%%%%%%%%%%%%%%%%%%%%%%%%%

Regression bootstrapping may resample:

\[
\boxed{
\text{observation pairs},
}
\]

or under suitable assumptions,

\[
\boxed{
\text{residuals}.
}
\]

The appropriate scheme depends on the regression design and error
structure.

Resampling must preserve the data-generating assumptions relevant to the
inferential target.

%%%%%%%%%%%%%%%%%%%%%%%%%%%%%%%%%%%%%%%%%%%%%%%%%%%%%%%%%%%%
\subsection{Bootstrap and Dependence}
\label{subsec:bootstrap-dependence}
%%%%%%%%%%%%%%%%%%%%%%%%%%%%%%%%%%%%%%%%%%%%%%%%%%%%%%%%%%%%

Ordinary observation-level bootstrap assumes independent sampling.

For time series, spatial data, clustered data, or repeated measurements,
independent resampling destroys dependence.

Specialized schemes include:

\[
\boxed{
\text{block bootstrap},
}
\]

\[
\boxed{
\text{cluster bootstrap},
}
\]

or model-based resampling.

%%%%%%%%%%%%%%%%%%%%%%%%%%%%%%%%%%%%%%%%%%%%%%%%%%%%%%%%%%%%
\subsection{Block Bootstrap Preview}
\label{subsec:block-bootstrap-preview}
%%%%%%%%%%%%%%%%%%%%%%%%%%%%%%%%%%%%%%%%%%%%%%%%%%%%%%%%%%%%

For dependent sequential data, blocks of consecutive observations are
resampled rather than individual observations.

This preserves short-range dependence within blocks.

The block length creates a bias--variance and dependence-preservation
tradeoff.

%%%%%%%%%%%%%%%%%%%%%%%%%%%%%%%%%%%%%%%%%%%%%%%%%%%%%%%%%%%%
\subsection{Jackknife}
\label{subsec:jackknife}
%%%%%%%%%%%%%%%%%%%%%%%%%%%%%%%%%%%%%%%%%%%%%%%%%%%%%%%%%%%%

The jackknife systematically removes one observation at a time.

Let

\[
\widehat{\theta}_{(-i)}
\]

be the estimate computed without observation \(i\).

There are

\[
\boxed{
n
}
\]

leave-one-out estimates.

%%%%%%%%%%%%%%%%%%%%%%%%%%%%%%%%%%%%%%%%%%%%%%%%%%%%%%%%%%%%
\subsection{Jackknife Mean}
\label{subsec:jackknife-mean}
%%%%%%%%%%%%%%%%%%%%%%%%%%%%%%%%%%%%%%%%%%%%%%%%%%%%%%%%%%%%

Define

\[
\boxed{
\overline{\theta}_{(\cdot)}
=
\frac1n
\sum_{i=1}^{n}
\widehat{\theta}_{(-i)}.
}
\]

A jackknife bias estimate is

\[
\boxed{
\widehat{\operatorname{Bias}}_{jack}
=
(n-1)
\left(
\overline{\theta}_{(\cdot)}
-
\widehat{\theta}
\right).
}
\]

%%%%%%%%%%%%%%%%%%%%%%%%%%%%%%%%%%%%%%%%%%%%%%%%%%%%%%%%%%%%
\subsection{Jackknife Standard Error}
\label{subsec:jackknife-standard-error}
%%%%%%%%%%%%%%%%%%%%%%%%%%%%%%%%%%%%%%%%%%%%%%%%%%%%%%%%%%%%

A jackknife standard-error estimate is

\[
\boxed{
\widehat{SE}_{jack}
=
\sqrt{
\frac{
n-1
}{
n
}
\sum_{i=1}^{n}
\left(
\widehat{\theta}_{(-i)}
-
\overline{\theta}_{(\cdot)}
\right)^2
}.
}
\]

Jackknife methods work especially well for sufficiently smooth
statistics.

%%%%%%%%%%%%%%%%%%%%%%%%%%%%%%%%%%%%%%%%%%%%%%%%%%%%%%%%%%%%
\subsection{Permutation Tests}
\label{subsec:permutation-tests-computational}
%%%%%%%%%%%%%%%%%%%%%%%%%%%%%%%%%%%%%%%%%%%%%%%%%%%%%%%%%%%%

Permutation tests construct a null distribution by rearranging labels or
assignments in a manner justified by the null hypothesis.

Suppose a statistic is

\[
\boxed{
T_{obs}.
}
\]

For each valid permutation \(\pi\), compute

\[
\boxed{
T_\pi.
}
\]

The permutation distribution provides the null reference distribution.

%%%%%%%%%%%%%%%%%%%%%%%%%%%%%%%%%%%%%%%%%%%%%%%%%%%%%%%%%%%%
\subsection{Exact Permutation Test}
\label{subsec:exact-permutation-test}
%%%%%%%%%%%%%%%%%%%%%%%%%%%%%%%%%%%%%%%%%%%%%%%%%%%%%%%%%%%%

If all valid permutations can be enumerated, the exact permutation
\(p\)-value for a right-tailed test is

\[
\boxed{
p
=
\frac{
\#\{
T_\pi\geq T_{obs}
\}
}{
\#\{
\text{valid permutations}
\}
}.
}
\]

The permutation scheme must reflect the null exchangeability structure.

%%%%%%%%%%%%%%%%%%%%%%%%%%%%%%%%%%%%%%%%%%%%%%%%%%%%%%%%%%%%
\subsection{Monte Carlo Permutation Test}
\label{subsec:monte-carlo-permutation-test}
%%%%%%%%%%%%%%%%%%%%%%%%%%%%%%%%%%%%%%%%%%%%%%%%%%%%%%%%%%%%

When the number of permutations is too large, sample

\[
B
\]

random valid permutations.

If

\[
B_{extreme}
\]

sampled statistics are at least as extreme as the observed statistic, a
common finite-simulation correction is

\[
\boxed{
p
=
\frac{
B_{extreme}+1
}{
B+1
}.
}
\]

The plus-one correction avoids reporting an impossible Monte Carlo
\(p\)-value of exactly zero.

%%%%%%%%%%%%%%%%%%%%%%%%%%%%%%%%%%%%%%%%%%%%%%%%%%%%%%%%%%%%
\subsection{Randomization Tests}
\label{subsec:randomization-tests-computational}
%%%%%%%%%%%%%%%%%%%%%%%%%%%%%%%%%%%%%%%%%%%%%%%%%%%%%%%%%%%%

In randomized experiments, the treatment-assignment mechanism defines a
natural reference distribution.

Under a sharp null hypothesis, outcomes are fixed and treatment labels
are reassigned according to the original randomization scheme.

This produces a design-based randomization distribution.

%%%%%%%%%%%%%%%%%%%%%%%%%%%%%%%%%%%%%%%%%%%%%%%%%%%%%%%%%%%%
\subsection{Permutation Versus Bootstrap}
\label{subsec:permutation-versus-bootstrap}
%%%%%%%%%%%%%%%%%%%%%%%%%%%%%%%%%%%%%%%%%%%%%%%%%%%%%%%%%%%%

Bootstrap methods approximate sampling distributions by resampling data
from an estimated population.

Permutation methods approximate null distributions by rearranging labels
or assignments under a null invariance principle.

Thus:

\[
\boxed{
\text{bootstrap}
\rightarrow
\text{sampling uncertainty},
}
\]

while

\[
\boxed{
\text{permutation}
\rightarrow
\text{null reference distribution}.
}
\]

%%%%%%%%%%%%%%%%%%%%%%%%%%%%%%%%%%%%%%%%%%%%%%%%%%%%%%%%%%%%
\subsection{Markov Chain Monte Carlo}
\label{subsec:mcmc-computational-statistics}
%%%%%%%%%%%%%%%%%%%%%%%%%%%%%%%%%%%%%%%%%%%%%%%%%%%%%%%%%%%%

Markov chain Monte Carlo constructs a dependent sequence

\[
\boxed{
\theta^{(1)},
\theta^{(2)},
\ldots
}
\]

whose stationary distribution is a target distribution

\[
\pi(\theta).
\]

In Bayesian statistics,

\[
\pi(\theta)
\]

is usually the posterior distribution.

%%%%%%%%%%%%%%%%%%%%%%%%%%%%%%%%%%%%%%%%%%%%%%%%%%%%%%%%%%%%
\subsection{Why MCMC Is Useful}
\label{subsec:why-mcmc-useful}
%%%%%%%%%%%%%%%%%%%%%%%%%%%%%%%%%%%%%%%%%%%%%%%%%%%%%%%%%%%%

Suppose the posterior is

\[
\pi(\theta\mid x)
=
\frac{
h(\theta)
}{
C
},
\]

where the normalizing constant \(C\) is unknown.

Many MCMC methods require only ratios of posterior densities.

Then \(C\) cancels.

This makes MCMC useful for high-dimensional intractable posteriors.

%%%%%%%%%%%%%%%%%%%%%%%%%%%%%%%%%%%%%%%%%%%%%%%%%%%%%%%%%%%%
\subsection{Metropolis--Hastings Algorithm}
\label{subsec:metropolis-hastings-computational}
%%%%%%%%%%%%%%%%%%%%%%%%%%%%%%%%%%%%%%%%%%%%%%%%%%%%%%%%%%%%

At current state

\[
\theta^{(m)},
\]

propose

\[
\theta'
\sim
q
(
\theta'\mid\theta^{(m)}
).
\]

Accept \(\theta'\) with probability

\[
\boxed{
\alpha
=
\min
\left\{
1,
\frac{
\pi(\theta')
q
(
\theta^{(m)}\mid\theta'
)
}{
\pi
(
\theta^{(m)}
)
q
(
\theta'\mid\theta^{(m)}
)
}
\right\}.
}
\]

If rejected,

\[
\boxed{
\theta^{(m+1)}
=
\theta^{(m)}.
}
\]

%%%%%%%%%%%%%%%%%%%%%%%%%%%%%%%%%%%%%%%%%%%%%%%%%%%%%%%%%%%%
\subsection{Symmetric Proposal}
\label{subsec:symmetric-mh-proposal}
%%%%%%%%%%%%%%%%%%%%%%%%%%%%%%%%%%%%%%%%%%%%%%%%%%%%%%%%%%%%

If

\[
q
(
\theta'\mid\theta
)
=
q
(
\theta\mid\theta'
),
\]

the proposal terms cancel.

Then

\[
\boxed{
\alpha
=
\min
\left\{
1,
\frac{
\pi(\theta')
}{
\pi(\theta)
}
\right\}.
}
\]

Random-walk proposals are common examples.

%%%%%%%%%%%%%%%%%%%%%%%%%%%%%%%%%%%%%%%%%%%%%%%%%%%%%%%%%%%%
\subsection{Detailed Balance}
\label{subsec:detailed-balance}
%%%%%%%%%%%%%%%%%%%%%%%%%%%%%%%%%%%%%%%%%%%%%%%%%%%%%%%%%%%%

A transition kernel satisfies detailed balance with respect to \(\pi\) if

\[
\boxed{
\pi(x)
P(x,dy)
=
\pi(y)
P(y,dx).
}
\]

Detailed balance implies that \(\pi\) is stationary.

Metropolis--Hastings is constructed to satisfy this relationship under
appropriate conditions.

%%%%%%%%%%%%%%%%%%%%%%%%%%%%%%%%%%%%%%%%%%%%%%%%%%%%%%%%%%%%
\subsection{Gibbs Sampling}
\label{subsec:gibbs-sampling-computational}
%%%%%%%%%%%%%%%%%%%%%%%%%%%%%%%%%%%%%%%%%%%%%%%%%%%%%%%%%%%%

Suppose

\[
\boldsymbol{\theta}
=
(
\theta_1,\ldots,\theta_p
).
\]

Gibbs sampling repeatedly draws each component from its full conditional:

\[
\boxed{
\theta_j
\sim
\pi
(
\theta_j
\mid
\boldsymbol{\theta}_{-j},
x
).
}
\]

If all full conditionals are easy to sample from, Gibbs sampling can be
simple to implement.

%%%%%%%%%%%%%%%%%%%%%%%%%%%%%%%%%%%%%%%%%%%%%%%%%%%%%%%%%%%%
\subsection{MCMC Burn-In and Warmup}
\label{subsec:mcmc-burnin-computational}
%%%%%%%%%%%%%%%%%%%%%%%%%%%%%%%%%%%%%%%%%%%%%%%%%%%%%%%%%%%%

Early iterations may depend strongly on initialization or be used for
algorithm adaptation.

These are often treated as warmup iterations.

Discarding warmup does not guarantee convergence.

The post-warmup chain must still mix adequately.

%%%%%%%%%%%%%%%%%%%%%%%%%%%%%%%%%%%%%%%%%%%%%%%%%%%%%%%%%%%%
\subsection{MCMC Autocorrelation}
\label{subsec:mcmc-autocorrelation-computational}
%%%%%%%%%%%%%%%%%%%%%%%%%%%%%%%%%%%%%%%%%%%%%%%%%%%%%%%%%%%%

MCMC draws are dependent.

Let

\[
\rho_k
\]

denote autocorrelation at lag \(k\) in the simulated chain.

Strong positive autocorrelation means the chain explores the target
slowly.

This reduces effective information.

%%%%%%%%%%%%%%%%%%%%%%%%%%%%%%%%%%%%%%%%%%%%%%%%%%%%%%%%%%%%
\subsection{Integrated Autocorrelation Time}
\label{subsec:integrated-autocorrelation-time}
%%%%%%%%%%%%%%%%%%%%%%%%%%%%%%%%%%%%%%%%%%%%%%%%%%%%%%%%%%%%

A conceptual integrated autocorrelation time is

\[
\boxed{
\tau
=
1
+
2
\sum_{k=1}^{\infty}
\rho_k,
}
\]

when the series is well-defined.

Then approximately

\[
\boxed{
ESS
\approx
\frac{
M
}{
\tau
}.
}
\]

Thus highly correlated chains have smaller effective sample sizes.

%%%%%%%%%%%%%%%%%%%%%%%%%%%%%%%%%%%%%%%%%%%%%%%%%%%%%%%%%%%%
\subsection{MCMC Monte Carlo Standard Error}
\label{subsec:mcmc-mcse}
%%%%%%%%%%%%%%%%%%%%%%%%%%%%%%%%%%%%%%%%%%%%%%%%%%%%%%%%%%%%

If posterior quantity \(g(\theta)\) has variance

\[
\sigma_g^2,
\]

a rough MCMC standard error is

\[
\boxed{
MCSE
\approx
\sqrt{
\frac{
\sigma_g^2
}{
ESS
}
}.
}
\]

Thus increasing raw iterations is useful only if the chain produces
additional effective information.

%%%%%%%%%%%%%%%%%%%%%%%%%%%%%%%%%%%%%%%%%%%%%%%%%%%%%%%%%%%%
\subsection{MCMC Diagnostics}
\label{subsec:mcmc-diagnostics-computational}
%%%%%%%%%%%%%%%%%%%%%%%%%%%%%%%%%%%%%%%%%%%%%%%%%%%%%%%%%%%%

Useful diagnostics include:

\[
\boxed{
\text{trace plots},
}
\]

\[
\boxed{
\text{autocorrelation plots},
}
\]

\[
\boxed{
\text{effective sample sizes},
}
\]

\[
\boxed{
\widehat{R},
}
\]

and comparisons across multiple chains.

No single diagnostic proves convergence.

%%%%%%%%%%%%%%%%%%%%%%%%%%%%%%%%%%%%%%%%%%%%%%%%%%%%%%%%%%%%
\subsection{Sequential Monte Carlo Preview}
\label{subsec:sequential-monte-carlo-preview}
%%%%%%%%%%%%%%%%%%%%%%%%%%%%%%%%%%%%%%%%%%%%%%%%%%%%%%%%%%%%

Sequential Monte Carlo represents a sequence of target distributions
using weighted particles.

At each stage, particles may undergo:

\[
\boxed{
\text{weighting},
}
\]

\[
\boxed{
\text{resampling},
}
\]

and

\[
\boxed{
\text{propagation}.
}
\]

These methods are widely used in state-space and dynamic models.

%%%%%%%%%%%%%%%%%%%%%%%%%%%%%%%%%%%%%%%%%%%%%%%%%%%%%%%%%%%%
\subsection{Particle Filtering Preview}
\label{subsec:particle-filtering-preview}
%%%%%%%%%%%%%%%%%%%%%%%%%%%%%%%%%%%%%%%%%%%%%%%%%%%%%%%%%%%%

A particle filter approximates the filtering distribution

\[
\boxed{
p
(
x_t
\mid
y_1,\ldots,y_t
)
}
\]

using a weighted collection of simulated states.

Unlike the Kalman filter, particle methods can handle nonlinear and
non-Gaussian state-space models.

%%%%%%%%%%%%%%%%%%%%%%%%%%%%%%%%%%%%%%%%%%%%%%%%%%%%%%%%%%%%
\subsection{Expectation--Maximization Algorithm}
\label{subsec:em-algorithm}
%%%%%%%%%%%%%%%%%%%%%%%%%%%%%%%%%%%%%%%%%%%%%%%%%%%%%%%%%%%%

The expectation--maximization, or EM, algorithm is used when a model
contains latent or missing data.

Let

\[
Z
\]

denote latent variables and

\[
X
\]

the observed data.

The complete-data log-likelihood is

\[
\boxed{
\ell_c
(
\theta;X,Z
).
}
\]

%%%%%%%%%%%%%%%%%%%%%%%%%%%%%%%%%%%%%%%%%%%%%%%%%%%%%%%%%%%%
\subsection{E-Step}
\label{subsec:em-e-step}
%%%%%%%%%%%%%%%%%%%%%%%%%%%%%%%%%%%%%%%%%%%%%%%%%%%%%%%%%%%%

Given current parameter estimate

\[
\theta^{(k)},
\]

the E-step calculates

\[
\boxed{
Q
(
\theta\mid\theta^{(k)}
)
=
\mathbb{E}
\left[
\ell_c
(
\theta;X,Z
)
\mid
X,\theta^{(k)}
\right].
}
\]

This replaces missing or latent information by its conditional expected
log-likelihood contribution.

%%%%%%%%%%%%%%%%%%%%%%%%%%%%%%%%%%%%%%%%%%%%%%%%%%%%%%%%%%%%
\subsection{M-Step}
\label{subsec:em-m-step}
%%%%%%%%%%%%%%%%%%%%%%%%%%%%%%%%%%%%%%%%%%%%%%%%%%%%%%%%%%%%

The M-step sets

\[
\boxed{
\theta^{(k+1)}
=
\arg\max_\theta
Q
(
\theta\mid\theta^{(k)}
).
}
\]

Thus EM alternates:

\[
\boxed{
\text{Expectation}
\rightarrow
\text{Maximization}.
}
\]

%%%%%%%%%%%%%%%%%%%%%%%%%%%%%%%%%%%%%%%%%%%%%%%%%%%%%%%%%%%%
\subsection{EM Monotonicity}
\label{subsec:em-monotonicity}
%%%%%%%%%%%%%%%%%%%%%%%%%%%%%%%%%%%%%%%%%%%%%%%%%%%%%%%%%%%%

Under standard conditions, each EM iteration does not decrease the
observed-data likelihood:

\[
\boxed{
L
(
\theta^{(k+1)}
)
\geq
L
(
\theta^{(k)}
).
}
\]

However, EM can converge slowly and may converge to a local optimum.

%%%%%%%%%%%%%%%%%%%%%%%%%%%%%%%%%%%%%%%%%%%%%%%%%%%%%%%%%%%%
\subsection{Mixture-Model Example}
\label{subsec:em-mixture-model}
%%%%%%%%%%%%%%%%%%%%%%%%%%%%%%%%%%%%%%%%%%%%%%%%%%%%%%%%%%%%

Suppose

\[
X_i
\]

arises from one of \(K\) latent mixture components.

Introduce latent indicator

\[
\boxed{
Z_i
\in
\{1,\ldots,K\}.
}
\]

The E-step estimates posterior probabilities of component membership.

The M-step updates component parameters using those probabilities as
weights.

%%%%%%%%%%%%%%%%%%%%%%%%%%%%%%%%%%%%%%%%%%%%%%%%%%%%%%%%%%%%
\subsection{Simulation Studies}
\label{subsec:simulation-studies}
%%%%%%%%%%%%%%%%%%%%%%%%%%%%%%%%%%%%%%%%%%%%%%%%%%%%%%%%%%%%

A simulation study evaluates statistical procedures under controlled
data-generating conditions.

Typical goals include estimating:

\[
\boxed{
\text{bias},
}
\]

\[
\boxed{
\text{variance},
}
\]

\[
\boxed{
\text{MSE},
}
\]

\[
\boxed{
\text{coverage},
}
\]

\[
\boxed{
\text{Type I error},
}
\]

or

\[
\boxed{
\text{power}.
}
\]

%%%%%%%%%%%%%%%%%%%%%%%%%%%%%%%%%%%%%%%%%%%%%%%%%%%%%%%%%%%%
\subsection{Basic Simulation Study Design}
\label{subsec:simulation-study-design}
%%%%%%%%%%%%%%%%%%%%%%%%%%%%%%%%%%%%%%%%%%%%%%%%%%%%%%%%%%%%

A simulation study may proceed as follows:

\begin{enumerate}
    \item choose known population parameters;
    \item generate a synthetic dataset;
    \item apply the statistical procedure;
    \item save the estimate or test result;
    \item repeat many times;
    \item summarize the repeated results.
\end{enumerate}

Because the true parameter is known, estimator properties can be measured
directly.

%%%%%%%%%%%%%%%%%%%%%%%%%%%%%%%%%%%%%%%%%%%%%%%%%%%%%%%%%%%%
\subsection{Monte Carlo Bias Estimate}
\label{subsec:simulation-bias}
%%%%%%%%%%%%%%%%%%%%%%%%%%%%%%%%%%%%%%%%%%%%%%%%%%%%%%%%%%%%

Suppose simulation replication \(r\) produces estimator

\[
\widehat{\theta}^{(r)}.
\]

Across \(R\) simulations,

\[
\boxed{
\widehat{\operatorname{Bias}}
=
\frac1R
\sum_{r=1}^{R}
\widehat{\theta}^{(r)}
-
\theta.
}
\]

%%%%%%%%%%%%%%%%%%%%%%%%%%%%%%%%%%%%%%%%%%%%%%%%%%%%%%%%%%%%
\subsection{Monte Carlo MSE Estimate}
\label{subsec:simulation-mse}
%%%%%%%%%%%%%%%%%%%%%%%%%%%%%%%%%%%%%%%%%%%%%%%%%%%%%%%%%%%%

The simulation estimate of MSE is

\[
\boxed{
\widehat{MSE}
=
\frac1R
\sum_{r=1}^{R}
\left(
\widehat{\theta}^{(r)}
-
\theta
\right)^2.
}
\]

This approximates the true repeated-sampling mean squared error.

%%%%%%%%%%%%%%%%%%%%%%%%%%%%%%%%%%%%%%%%%%%%%%%%%%%%%%%%%%%%
\subsection{Coverage Simulation}
\label{subsec:coverage-simulation}
%%%%%%%%%%%%%%%%%%%%%%%%%%%%%%%%%%%%%%%%%%%%%%%%%%%%%%%%%%%%

Suppose interval

\[
C_r
\]

is constructed in simulation \(r\).

The empirical coverage is

\[
\boxed{
\widehat{Coverage}
=
\frac1R
\sum_{r=1}^{R}
\mathbf{1}
\{
\theta\in C_r
\}.
}
\]

A nominal \(95\%\) interval should ideally yield coverage near

\[
\boxed{
0.95.
}
\]

%%%%%%%%%%%%%%%%%%%%%%%%%%%%%%%%%%%%%%%%%%%%%%%%%%%%%%%%%%%%
\subsection{Simulating Type I Error}
\label{subsec:simulation-type-i}
%%%%%%%%%%%%%%%%%%%%%%%%%%%%%%%%%%%%%%%%%%%%%%%%%%%%%%%%%%%%

Generate data under the null hypothesis.

Apply the test and record whether it rejects.

Then

\[
\boxed{
\widehat{\alpha}
=
\frac{
\text{number of rejections}
}{
R
}.
}
\]

This evaluates whether the procedure maintains its nominal significance
level.

%%%%%%%%%%%%%%%%%%%%%%%%%%%%%%%%%%%%%%%%%%%%%%%%%%%%%%%%%%%%
\subsection{Simulating Power}
\label{subsec:simulation-power}
%%%%%%%%%%%%%%%%%%%%%%%%%%%%%%%%%%%%%%%%%%%%%%%%%%%%%%%%%%%%

Generate data under a specified alternative.

Then

\[
\boxed{
\widehat{Power}
=
\frac{
\text{number of rejections}
}{
R
}.
}
\]

Repeating this for several effect sizes produces an empirical power curve.

%%%%%%%%%%%%%%%%%%%%%%%%%%%%%%%%%%%%%%%%%%%%%%%%%%%%%%%%%%%%
\subsection{Monte Carlo Error in Simulation Studies}
\label{subsec:simulation-monte-carlo-error}
%%%%%%%%%%%%%%%%%%%%%%%%%%%%%%%%%%%%%%%%%%%%%%%%%%%%%%%%%%%%

Simulation summaries themselves contain Monte Carlo error.

For a simulated rejection proportion

\[
\widehat{p},
\]

the approximate Monte Carlo standard error is

\[
\boxed{
\sqrt{
\frac{
\widehat{p}
(
1-\widehat{p}
)
}{
R
}
}.
}
\]

Thus reported simulation results should use enough replications for the
desired precision.

%%%%%%%%%%%%%%%%%%%%%%%%%%%%%%%%%%%%%%%%%%%%%%%%%%%%%%%%%%%%
\subsection{Reproducible Simulation}
\label{subsec:reproducible-simulation}
%%%%%%%%%%%%%%%%%%%%%%%%%%%%%%%%%%%%%%%%%%%%%%%%%%%%%%%%%%%%

A reproducible simulation should document:

\[
\boxed{
\text{data-generating mechanism},
}
\]

\[
\boxed{
\text{parameter values},
}
\]

\[
\boxed{
\text{sample sizes},
}
\]

\[
\boxed{
\text{number of replications},
}
\]

\[
\boxed{
\text{random-number settings},
}
\]

and

\[
\boxed{
\text{analysis procedure}.
}
\]

%%%%%%%%%%%%%%%%%%%%%%%%%%%%%%%%%%%%%%%%%%%%%%%%%%%%%%%%%%%%
\subsection{Parallel Computation}
\label{subsec:parallel-computation-statistics}
%%%%%%%%%%%%%%%%%%%%%%%%%%%%%%%%%%%%%%%%%%%%%%%%%%%%%%%%%%%%

Many simulation and resampling tasks contain independent repetitions.

For example:

\[
\boxed{
\text{bootstrap replications},
}
\]

\[
\boxed{
\text{permutation samples},
}
\]

and

\[
\boxed{
\text{simulation-study replications}.
}
\]

These can often be distributed across multiple processors.

%%%%%%%%%%%%%%%%%%%%%%%%%%%%%%%%%%%%%%%%%%%%%%%%%%%%%%%%%%%%
\subsection{Parallel Random Number Generation}
\label{subsec:parallel-random-number-generation}
%%%%%%%%%%%%%%%%%%%%%%%%%%%%%%%%%%%%%%%%%%%%%%%%%%%%%%%%%%%%

Parallel simulation requires care with pseudo-random streams.

Simply reusing the same seed on every worker can generate identical
sequences.

Reliable parallel random-number methods use independent or appropriately
separated streams.

%%%%%%%%%%%%%%%%%%%%%%%%%%%%%%%%%%%%%%%%%%%%%%%%%%%%%%%%%%%%
\subsection{Vectorization}
\label{subsec:vectorization-statistics}
%%%%%%%%%%%%%%%%%%%%%%%%%%%%%%%%%%%%%%%%%%%%%%%%%%%%%%%%%%%%

Vectorized computation performs operations on entire arrays or matrices
rather than using explicit scalar loops.

For example,

\[
\boxed{
X^T\mathbf{y}
}
\]

is a matrix operation.

Vectorization can improve both clarity and speed when supported by
optimized numerical libraries.

%%%%%%%%%%%%%%%%%%%%%%%%%%%%%%%%%%%%%%%%%%%%%%%%%%%%%%%%%%%%
\subsection{Sparse Matrices}
\label{subsec:sparse-matrices-statistics}
%%%%%%%%%%%%%%%%%%%%%%%%%%%%%%%%%%%%%%%%%%%%%%%%%%%%%%%%%%%%

A sparse matrix contains mostly zero entries.

Large models involving:

\[
\boxed{
\text{networks},
}
\]

\[
\boxed{
\text{spatial adjacency},
}
\]

or

\[
\boxed{
\text{high-dimensional design matrices}
}
\]

may benefit greatly from sparse storage and algorithms.

Storing every zero explicitly can waste enormous memory.

%%%%%%%%%%%%%%%%%%%%%%%%%%%%%%%%%%%%%%%%%%%%%%%%%%%%%%%%%%%%
\subsection{Computational Complexity}
\label{subsec:computational-complexity-statistics}
%%%%%%%%%%%%%%%%%%%%%%%%%%%%%%%%%%%%%%%%%%%%%%%%%%%%%%%%%%%%

Algorithms can be compared by how computational cost grows with problem
size.

For example, dense matrix inversion is often roughly

\[
\boxed{
O(p^3).
}
\]

Thus doubling \(p\) can increase computational cost by approximately a
factor of

\[
\boxed{
8.
}
\]

Complexity becomes critical in high-dimensional problems.

%%%%%%%%%%%%%%%%%%%%%%%%%%%%%%%%%%%%%%%%%%%%%%%%%%%%%%%%%%%%
\subsection{Memory Complexity}
\label{subsec:memory-complexity-statistics}
%%%%%%%%%%%%%%%%%%%%%%%%%%%%%%%%%%%%%%%%%%%%%%%%%%%%%%%%%%%%

A dense \(p\times p\) matrix requires storage proportional to

\[
\boxed{
p^2.
}
\]

For very large \(p\), covariance and kernel matrices can become
impractical even before arithmetic becomes too slow.

Algorithms may therefore rely on:

\[
\boxed{
\text{sparsity},
}
\]

\[
\boxed{
\text{low-rank approximations},
}
\]

or

\[
\boxed{
\text{iterative matrix methods}.
}
\]

%%%%%%%%%%%%%%%%%%%%%%%%%%%%%%%%%%%%%%%%%%%%%%%%%%%%%%%%%%%%
\subsection{Numerical Stability in Least Squares}
\label{subsec:numerical-stability-least-squares}
%%%%%%%%%%%%%%%%%%%%%%%%%%%%%%%%%%%%%%%%%%%%%%%%%%%%%%%%%%%%

The normal-equation estimator is

\[
\widehat{\boldsymbol{\beta}}
=
(
X^TX
)^{-1}
X^T\mathbf{y}.
\]

But forming

\[
X^TX
\]

can worsen conditioning.

QR or SVD-based least-squares solutions are usually more stable.

Thus mathematically equivalent formulas need not be computationally
equivalent.

%%%%%%%%%%%%%%%%%%%%%%%%%%%%%%%%%%%%%%%%%%%%%%%%%%%%%%%%%%%%
\subsection{Iterative Solvers}
\label{subsec:iterative-linear-solvers-statistics}
%%%%%%%%%%%%%%%%%%%%%%%%%%%%%%%%%%%%%%%%%%%%%%%%%%%%%%%%%%%%

For extremely large systems, direct matrix decompositions may be too
expensive.

Iterative solvers approximate the solution to

\[
\boxed{
A\mathbf{x}
=
\mathbf{b}
}
\]

through successive updates.

Such methods can exploit matrix sparsity and avoid storing dense
factorizations.

%%%%%%%%%%%%%%%%%%%%%%%%%%%%%%%%%%%%%%%%%%%%%%%%%%%%%%%%%%%%
\subsection{Numerical Integration}
\label{subsec:numerical-integration-statistics}
%%%%%%%%%%%%%%%%%%%%%%%%%%%%%%%%%%%%%%%%%%%%%%%%%%%%%%%%%%%%

Statistical expectations often have the form

\[
\boxed{
\int
g(\theta)
p(\theta)
\,d\theta.
}
\]

Low-dimensional integrals may be approximated using deterministic
quadrature.

High-dimensional integrals often require Monte Carlo or other specialized
methods.

%%%%%%%%%%%%%%%%%%%%%%%%%%%%%%%%%%%%%%%%%%%%%%%%%%%%%%%%%%%%
\subsection{Quadrature Preview}
\label{subsec:quadrature-preview}
%%%%%%%%%%%%%%%%%%%%%%%%%%%%%%%%%%%%%%%%%%%%%%%%%%%%%%%%%%%%

A numerical quadrature rule approximates

\[
\int_a^b
f(x)\,dx
\]

using weighted evaluations:

\[
\boxed{
\int_a^b
f(x)\,dx
\approx
\sum_{j=1}^{m}
w_j
f(x_j).
}
\]

Examples include trapezoidal, Simpson-type, and Gaussian quadrature
methods.

%%%%%%%%%%%%%%%%%%%%%%%%%%%%%%%%%%%%%%%%%%%%%%%%%%%%%%%%%%%%
\subsection{Curse of Dimensionality in Integration}
\label{subsec:integration-curse-dimensionality}
%%%%%%%%%%%%%%%%%%%%%%%%%%%%%%%%%%%%%%%%%%%%%%%%%%%%%%%%%%%%

A deterministic grid using \(m\) points per dimension requires

\[
\boxed{
m^d
}
\]

evaluations in \(d\) dimensions.

Thus quadrature becomes rapidly infeasible as dimension grows.

Monte Carlo methods are attractive because their basic convergence rate
does not directly deteriorate exponentially with dimension in this way.

%%%%%%%%%%%%%%%%%%%%%%%%%%%%%%%%%%%%%%%%%%%%%%%%%%%%%%%%%%%%
\subsection{Optimization in Statistical Learning}
\label{subsec:optimization-statistical-learning}
%%%%%%%%%%%%%%%%%%%%%%%%%%%%%%%%%%%%%%%%%%%%%%%%%%%%%%%%%%%%

Modern predictive models frequently solve penalized objectives such as

\[
\boxed{
\min_{\boldsymbol{\beta}}
\left[
L
(
\boldsymbol{\beta}
)
+
\lambda
P
(
\boldsymbol{\beta}
)
\right].
}
\]

Different penalties require different numerical algorithms.

Examples include:

\[
\boxed{
\text{gradient methods},
}
\]

\[
\boxed{
\text{coordinate descent},
}
\]

and

\[
\boxed{
\text{proximal algorithms}.
}
\]

%%%%%%%%%%%%%%%%%%%%%%%%%%%%%%%%%%%%%%%%%%%%%%%%%%%%%%%%%%%%
\subsection{Proximal Methods Preview}
\label{subsec:proximal-methods-preview}
%%%%%%%%%%%%%%%%%%%%%%%%%%%%%%%%%%%%%%%%%%%%%%%%%%%%%%%%%%%%

Suppose an objective is

\[
f(\theta)+g(\theta),
\]

where \(f\) is smooth and \(g\) may be nonsmooth.

A proximal-gradient update has form

\[
\boxed{
\theta^{(k+1)}
=
\operatorname{prox}_{\eta g}
\left[
\theta^{(k)}
-
\eta
\nabla f
(
\theta^{(k)}
)
\right].
}
\]

This framework is useful for \(\ell_1\)-penalized problems.

%%%%%%%%%%%%%%%%%%%%%%%%%%%%%%%%%%%%%%%%%%%%%%%%%%%%%%%%%%%%
\subsection{Reproducible Statistical Computation}
\label{subsec:reproducible-statistical-computation}
%%%%%%%%%%%%%%%%%%%%%%%%%%%%%%%%%%%%%%%%%%%%%%%%%%%%%%%%%%%%

A reproducible analysis should record:

\[
\boxed{
\text{data version},
}
\]

\[
\boxed{
\text{code},
}
\]

\[
\boxed{
\text{software versions},
}
\]

\[
\boxed{
\text{random seeds},
}
\]

\[
\boxed{
\text{model settings},
}
\]

and

\[
\boxed{
\text{preprocessing decisions}.
}
\]

Reproducibility is part of statistical validity.

%%%%%%%%%%%%%%%%%%%%%%%%%%%%%%%%%%%%%%%%%%%%%%%%%%%%%%%%%%%%
\subsection{Computational Diagnostics}
\label{subsec:computational-diagnostics}
%%%%%%%%%%%%%%%%%%%%%%%%%%%%%%%%%%%%%%%%%%%%%%%%%%%%%%%%%%%%

A statistical result should not be trusted solely because software
returned a number.

Check for:

\[
\boxed{
\text{optimizer convergence},
}
\]

\[
\boxed{
\text{gradient size},
}
\]

\[
\boxed{
\text{boundary estimates},
}
\]

\[
\boxed{
\text{singular matrices},
}
\]

\[
\boxed{
\text{MCMC mixing},
}
\]

and

\[
\boxed{
\text{simulation error}.
}
\]

%%%%%%%%%%%%%%%%%%%%%%%%%%%%%%%%%%%%%%%%%%%%%%%%%%%%%%%%%%%%
\subsection{Algorithm Verification}
\label{subsec:algorithm-verification}
%%%%%%%%%%%%%%%%%%%%%%%%%%%%%%%%%%%%%%%%%%%%%%%%%%%%%%%%%%%%

When possible, verify a computational method using a problem with a known
analytical answer.

For example, an MCMC algorithm can be tested on a simple normal target
whose moments are known.

Agreement does not prove correctness for all problems, but disagreement
reveals implementation errors.

%%%%%%%%%%%%%%%%%%%%%%%%%%%%%%%%%%%%%%%%%%%%%%%%%%%%%%%%%%%%
\subsection{Sensitivity to Initialization}
\label{subsec:sensitivity-initialization}
%%%%%%%%%%%%%%%%%%%%%%%%%%%%%%%%%%%%%%%%%%%%%%%%%%%%%%%%%%%%

Some algorithms depend strongly on starting values.

Examples include:

\[
\boxed{
\text{mixture-model EM},
}
\]

\[
\boxed{
k\text{-means},
}
\]

and

\[
\boxed{
\text{nonconvex likelihood optimization}.
}
\]

Multiple initializations can expose instability.

%%%%%%%%%%%%%%%%%%%%%%%%%%%%%%%%%%%%%%%%%%%%%%%%%%%%%%%%%%%%
\subsection{Sensitivity to Tolerance}
\label{subsec:sensitivity-tolerance}
%%%%%%%%%%%%%%%%%%%%%%%%%%%%%%%%%%%%%%%%%%%%%%%%%%%%%%%%%%%%

A convergence tolerance that is too loose may stop an algorithm before
adequate accuracy is reached.

An excessively strict tolerance can waste computation or fail because of
floating-point limitations.

Important analyses should be checked for stability across reasonable
tolerances.

%%%%%%%%%%%%%%%%%%%%%%%%%%%%%%%%%%%%%%%%%%%%%%%%%%%%%%%%%%%%
\subsection{High-Dimensional Statistical Computation}
\label{subsec:high-dimensional-statistical-computation}
%%%%%%%%%%%%%%%%%%%%%%%%%%%%%%%%%%%%%%%%%%%%%%%%%%%%%%%%%%%%

When

\[
p
\]

is extremely large, standard dense computations may become infeasible.

Possible strategies include:

\[
\boxed{
\text{regularization},
}
\]

\[
\boxed{
\text{sparse matrices},
}
\]

\[
\boxed{
\text{iterative solvers},
}
\]

\[
\boxed{
\text{low-rank approximations},
}
\]

and

\[
\boxed{
\text{distributed computation}.
}
\]

%%%%%%%%%%%%%%%%%%%%%%%%%%%%%%%%%%%%%%%%%%%%%%%%%%%%%%%%%%%%
\subsection{Low-Rank Approximation}
\label{subsec:low-rank-approximation-statistics}
%%%%%%%%%%%%%%%%%%%%%%%%%%%%%%%%%%%%%%%%%%%%%%%%%%%%%%%%%%%%

If a matrix is approximately low rank, one can retain only its largest
singular components.

For

\[
X
=
UDV^T,
\]

a rank-\(r\) approximation is

\[
\boxed{
X_r
=
U_r
D_r
V_r^T.
}
\]

This reduces storage and computation while retaining dominant structure.

%%%%%%%%%%%%%%%%%%%%%%%%%%%%%%%%%%%%%%%%%%%%%%%%%%%%%%%%%%%%
\subsection{Computational Statistical Workflow}
\label{subsec:computational-statistics-workflow}
%%%%%%%%%%%%%%%%%%%%%%%%%%%%%%%%%%%%%%%%%%%%%%%%%%%%%%%%%%%%

A practical computational workflow is:

\begin{enumerate}
    \item formulate the statistical target;
    \item determine whether an analytical solution exists;
    \item choose an appropriate numerical algorithm;
    \item choose a numerically stable parameterization;
    \item initialize the algorithm carefully;
    \item set convergence tolerances;
    \item monitor objective values or diagnostic quantities;
    \item verify convergence;
    \item compare multiple initializations when appropriate;
    \item quantify Monte Carlo error if simulation is used;
    \item verify the method on known special cases;
    \item assess sensitivity to algorithm settings;
    \item use stable matrix factorizations;
    \item preserve random seeds and software settings;
    \item document computational limitations.
\end{enumerate}

%%%%%%%%%%%%%%%%%%%%%%%%%%%%%%%%%%%%%%%%%%%%%%%%%%%%%%%%%%%%
\subsection{Choosing a Computational Method}
\label{subsec:choosing-computational-method}
%%%%%%%%%%%%%%%%%%%%%%%%%%%%%%%%%%%%%%%%%%%%%%%%%%%%%%%%%%%%

For smooth likelihood optimization:

\[
\boxed{
\text{Newton, quasi-Newton, or gradient methods}.
}
\]

For lasso-type penalization:

\[
\boxed{
\text{coordinate descent or proximal methods}.
}
\]

For uncertainty from repeated sampling:

\[
\boxed{
\text{bootstrap}.
}
\]

For null distributions under exchangeability:

\[
\boxed{
\text{permutation or randomization tests}.
}
\]

For ordinary integrals under known simulation distributions:

\[
\boxed{
\text{Monte Carlo}.
}
\]

For difficult posterior distributions:

\[
\boxed{
\text{MCMC}.
}
\]

For latent-variable likelihoods:

\[
\boxed{
\text{EM}.
}
\]

%%%%%%%%%%%%%%%%%%%%%%%%%%%%%%%%%%%%%%%%%%%%%%%%%%%%%%%%%%%%
\subsection{Common Errors}
\label{subsec:computational-statistics-common-errors}
%%%%%%%%%%%%%%%%%%%%%%%%%%%%%%%%%%%%%%%%%%%%%%%%%%%%%%%%%%%%

\subsubsection{Assuming Software Output Is Automatically Correct}

A numerical optimizer may stop at:

\[
\boxed{
\text{a local optimum},
}
\]

\[
\boxed{
\text{a boundary},
}
\]

or

\[
\boxed{
\text{a nonconverged point}.
}
\]

Diagnostic messages must be inspected.

%%%%%%%%%%%%%%%%%%%%%%%%%%%%%%%%%%%%%%%%%%%%%%%%%%%%%%%%%%%%
\subsubsection{Ignoring Numerical Precision}

Very large or very small numbers can cause overflow or underflow.

Likelihood calculations should often use logarithmic forms.

%%%%%%%%%%%%%%%%%%%%%%%%%%%%%%%%%%%%%%%%%%%%%%%%%%%%%%%%%%%%
\subsubsection{Explicitly Inverting Matrices Unnecessarily}

Solving

\[
A\mathbf{x}
=
\mathbf{b}
\]

directly is usually preferable to computing

\[
A^{-1}
\]

and multiplying.

%%%%%%%%%%%%%%%%%%%%%%%%%%%%%%%%%%%%%%%%%%%%%%%%%%%%%%%%%%%%
\subsubsection{Using Normal Equations in Ill-Conditioned Problems}

Forming

\[
X^TX
\]

can worsen conditioning.

QR or SVD methods are often more stable.

%%%%%%%%%%%%%%%%%%%%%%%%%%%%%%%%%%%%%%%%%%%%%%%%%%%%%%%%%%%%
\subsubsection{Ignoring Multiple Local Optima}

Mixture models, clustering, and nonconvex objectives can have several
local solutions.

One starting point may not be enough.

%%%%%%%%%%%%%%%%%%%%%%%%%%%%%%%%%%%%%%%%%%%%%%%%%%%%%%%%%%%%
\subsubsection{Stopping Optimization Too Early}

A small number of iterations does not guarantee convergence.

The gradient, parameter changes, and objective changes should be checked.

%%%%%%%%%%%%%%%%%%%%%%%%%%%%%%%%%%%%%%%%%%%%%%%%%%%%%%%%%%%%
\subsubsection{Reporting Monte Carlo Estimates Without Monte Carlo Error}

A simulation result such as

\[
0.0317
\]

has finite-simulation uncertainty.

Report or at least assess its Monte Carlo standard error.

%%%%%%%%%%%%%%%%%%%%%%%%%%%%%%%%%%%%%%%%%%%%%%%%%%%%%%%%%%%%
\subsubsection{Reporting a Monte Carlo \(p\)-Value of Zero}

If only \(B\) random permutations are sampled, use a finite-simulation
correction such as

\[
\boxed{
\frac{
B_{extreme}+1
}{
B+1
}.
}
\]

The true \(p\)-value has not been established as exactly zero.

%%%%%%%%%%%%%%%%%%%%%%%%%%%%%%%%%%%%%%%%%%%%%%%%%%%%%%%%%%%%
\subsubsection{Using Ordinary Bootstrap for Dependent Data}

Independent resampling destroys serial, spatial, or cluster dependence.

Use a resampling scheme appropriate to the sampling structure.

%%%%%%%%%%%%%%%%%%%%%%%%%%%%%%%%%%%%%%%%%%%%%%%%%%%%%%%%%%%%
\subsubsection{Assuming Bootstrap Always Works}

The bootstrap can fail for irregular estimators, boundary problems,
extreme-value statistics, or inappropriate resampling schemes.

Its validity depends on the statistic and data-generating process.

%%%%%%%%%%%%%%%%%%%%%%%%%%%%%%%%%%%%%%%%%%%%%%%%%%%%%%%%%%%%
\subsubsection{Ignoring Bootstrap Monte Carlo Error}

The estimated bootstrap quantiles themselves depend on finite \(B\).

Too few bootstrap samples can create unstable intervals.

%%%%%%%%%%%%%%%%%%%%%%%%%%%%%%%%%%%%%%%%%%%%%%%%%%%%%%%%%%%%
\subsubsection{Confusing Permutation and Bootstrap Logic}

Bootstrap resampling approximates sampling variation.

Permutation resampling constructs a null reference distribution.

They answer different inferential questions.

%%%%%%%%%%%%%%%%%%%%%%%%%%%%%%%%%%%%%%%%%%%%%%%%%%%%%%%%%%%%
\subsubsection{Assuming MCMC Draws Are Independent}

MCMC samples are generally autocorrelated.

Raw iteration count overstates effective information.

%%%%%%%%%%%%%%%%%%%%%%%%%%%%%%%%%%%%%%%%%%%%%%%%%%%%%%%%%%%%
\subsubsection{Using Only One MCMC Chain}

A single chain can remain trapped in one region without revealing the
problem.

Multiple dispersed chains provide stronger diagnostics.

%%%%%%%%%%%%%%%%%%%%%%%%%%%%%%%%%%%%%%%%%%%%%%%%%%%%%%%%%%%%
\subsubsection{Assuming Warmup Fixes Poor Mixing}

Discarding early iterations does not correct an algorithm that fails to
explore the target distribution.

%%%%%%%%%%%%%%%%%%%%%%%%%%%%%%%%%%%%%%%%%%%%%%%%%%%%%%%%%%%%
\subsubsection{Thinning Automatically}

Thinning often discards useful draws.

It should be used only for a justified computational reason.

%%%%%%%%%%%%%%%%%%%%%%%%%%%%%%%%%%%%%%%%%%%%%%%%%%%%%%%%%%%%
\subsubsection{Ignoring Importance Weight Degeneracy}

A small number of huge weights can make importance-sampling estimates
effectively depend on only a few simulations.

%%%%%%%%%%%%%%%%%%%%%%%%%%%%%%%%%%%%%%%%%%%%%%%%%%%%%%%%%%%%
\subsubsection{Using the Same Random Stream Improperly in Parallel Code}

Duplicated random streams can make apparently independent simulation
replications identical or correlated.

%%%%%%%%%%%%%%%%%%%%%%%%%%%%%%%%%%%%%%%%%%%%%%%%%%%%%%%%%%%%
\subsubsection{Comparing Algorithms Without Equal Computational Budgets}

One method may appear more accurate simply because it used much more
computation.

Runtime, memory, convergence quality, and statistical accuracy should be
considered together.

%%%%%%%%%%%%%%%%%%%%%%%%%%%%%%%%%%%%%%%%%%%%%%%%%%%%%%%%%%%%
\subsubsection{Ignoring Reproducibility}

Without code, software settings, preprocessing details, and random-number
information, a computational statistical result may be impossible to
verify.

%%%%%%%%%%%%%%%%%%%%%%%%%%%%%%%%%%%%%%%%%%%%%%%%%%%%%%%%%%%%
\subsection{Applications}
\label{subsec:computational-statistics-applications}
%%%%%%%%%%%%%%%%%%%%%%%%%%%%%%%%%%%%%%%%%%%%%%%%%%%%%%%%%%%%

\begin{applicationbox}

\textbf{Bayesian Inference.}
MCMC, importance sampling, quadrature, and optimization make complex
posterior distributions computationally accessible.

\medskip

\textbf{Statistical Learning.}
Regularized regression, boosting, kernel methods, and high-dimensional
models rely on numerical optimization.

\medskip

\textbf{Environmental Science.}
Bootstrap, Monte Carlo simulation, spatial computation, and stochastic
models quantify uncertainty in environmental data.

\medskip

\textbf{Engineering.}
Simulation studies evaluate system reliability, estimator performance,
and uncertainty under complicated failure processes.

\medskip

\textbf{Medicine.}
Computational methods fit survival models, hierarchical models,
longitudinal models, and high-dimensional predictive systems.

\medskip

\textbf{Finance.}
Monte Carlo methods estimate risk, option values, portfolio uncertainty,
and stochastic-model quantities.

\medskip

\textbf{Operations Research.}
Simulation and optimization are central to queueing, scheduling,
inventory, and stochastic decision problems.

\medskip

\textbf{Scientific Computing.}
Matrix decompositions, iterative methods, parallel simulation, and
numerical integration support large-scale quantitative research.

\end{applicationbox}

%%%%%%%%%%%%%%%%%%%%%%%%%%%%%%%%%%%%%%%%%%%%%%%%%%%%%%%%%%%%
\subsection{Exercises}
\label{subsec:computational-statistics-exercises}
%%%%%%%%%%%%%%%%%%%%%%%%%%%%%%%%%%%%%%%%%%%%%%%%%%%%%%%%%%%%

\begin{exercise}[1]
Define computational statistics.
\end{exercise}

\begin{exercise}[1]
Distinguish sampling error and computational error.
\end{exercise}

\begin{exercise}[1]
Define numerical instability.
\end{exercise}

\begin{exercise}[1]
Define Monte Carlo integration.
\end{exercise}

\begin{exercise}[1]
Define bootstrap resampling.
\end{exercise}

\begin{exercise}[1]
Define a permutation test.
\end{exercise}

\begin{exercise}[1]
Define Markov chain Monte Carlo.
\end{exercise}

\begin{exercise}[1]
Define the EM algorithm.
\end{exercise}

\begin{exercise}[1]
Define a random seed.
\end{exercise}

\begin{exercise}[1]
Define effective sample size.
\end{exercise}

\begin{exercise}[1]
Define importance sampling.
\end{exercise}

\begin{exercise}[1]
Define a condition number.
\end{exercise}

\begin{exercise}[2]
Suppose an ordinary Monte Carlo estimator has standard deviation \(8\)
per simulated draw and uses

\[
M=400.
\]

Compute its Monte Carlo standard error.
\end{exercise}

\begin{exercise}[2]
How many simulations are required to reduce Monte Carlo standard error by
one-half, approximately?
\end{exercise}

\begin{exercise}[2]
Suppose

\[
U=0.25
\]

and

\[
\lambda=2.
\]

Use inverse transform sampling to generate an exponential random variable
using

\[
X=-\frac1\lambda\log U.
\]
\end{exercise}

\begin{exercise}[2]
Suppose normalized importance weights are

\[
0.5,\ 0.25,\ 0.25.
\]

Compute the heuristic ESS

\[
\frac1{\sum w_i^2}.
\]
\end{exercise}

\begin{exercise}[2]
Suppose bootstrap estimates have standard deviation

\[
2.4.
\]

What quantity does this approximate?
\end{exercise}

\begin{exercise}[2]
Suppose \(1000\) random permutations are sampled and \(19\) produce
statistics at least as extreme as the observed statistic.

Compute the plus-one Monte Carlo permutation \(p\)-value.
\end{exercise}

\begin{exercise}[2]
Suppose an MCMC run contains \(10000\) draws with estimated integrated
autocorrelation time \(5\).

Approximate its effective sample size.
\end{exercise}

\begin{exercise}[2]
Suppose a dense matrix has dimension

\[
p=1000.
\]

Approximately how many entries must be stored?
\end{exercise}

\begin{exercise}[2]
If dense matrix inversion scales approximately as \(p^3\), by what factor
does computational cost increase when \(p\) doubles?
\end{exercise}

\begin{exercise}[3]
Explain why log-likelihood computation is often more numerically stable
than direct likelihood computation.
\end{exercise}

\begin{exercise}[3]
Derive the log-sum-exp identity.
\end{exercise}

\begin{exercise}[3]
Explain why subtracting nearly equal floating-point numbers can be
unstable.
\end{exercise}

\begin{exercise}[3]
Derive the gradient-descent update from the local direction of steepest
descent.
\end{exercise}

\begin{exercise}[3]
Derive Newton's method using a second-order Taylor expansion.
\end{exercise}

\begin{exercise}[3]
Compare Newton--Raphson and Fisher scoring.
\end{exercise}

\begin{exercise}[3]
Explain why reparameterizing \(\sigma=e^\eta\) removes a positivity
constraint.
\end{exercise}

\begin{exercise}[3]
Derive the soft-thresholding solution for a one-dimensional lasso
problem.
\end{exercise}

\begin{exercise}[3]
Explain why QR decomposition is numerically preferable to the normal
equations in ill-conditioned least-squares problems.
\end{exercise}

\begin{exercise}[3]
Explain how SVD identifies rank deficiency.
\end{exercise}

\begin{exercise}[3]
Explain the statistical meaning of a large matrix condition number.
\end{exercise}

\begin{exercise}[3]
Prove consistency of the basic Monte Carlo estimator using the law of
large numbers.
\end{exercise}

\begin{exercise}[3]
Derive the variance of the basic Monte Carlo mean.
\end{exercise}

\begin{exercise}[3]
Explain why Monte Carlo error decreases at rate \(M^{-1/2}\).
\end{exercise}

\begin{exercise}[3]
Derive inverse transform sampling.
\end{exercise}

\begin{exercise}[3]
Show that acceptance--rejection sampling produces draws from the target
density.
\end{exercise}

\begin{exercise}[3]
Derive the importance-sampling identity.
\end{exercise}

\begin{exercise}[3]
Explain why poor proposal-tail coverage can cause infinite or enormous
importance-sampling variance.
\end{exercise}

\begin{exercise}[3]
Derive the optimal control-variate coefficient.
\end{exercise}

\begin{exercise}[3]
Explain how antithetic variables can reduce variance.
\end{exercise}

\begin{exercise}[3]
Derive the bootstrap standard-error formula.
\end{exercise}

\begin{exercise}[3]
Compare percentile, basic, and normal bootstrap confidence intervals.
\end{exercise}

\begin{exercise}[3]
Explain the difference between parametric and nonparametric bootstrap.
\end{exercise}

\begin{exercise}[3]
Explain why an ordinary bootstrap is inappropriate for strongly
autocorrelated time-series data.
\end{exercise}

\begin{exercise}[3]
Derive the jackknife standard-error formula.
\end{exercise}

\begin{exercise}[3]
Explain the null exchangeability condition underlying permutation tests.
\end{exercise}

\begin{exercise}[3]
Explain why the plus-one permutation correction is used in Monte Carlo
permutation tests.
\end{exercise}

\begin{exercise}[3]
Explain the difference between exact and sampled permutation tests.
\end{exercise}

\begin{exercise}[3]
Explain why Metropolis--Hastings does not require the posterior
normalizing constant.
\end{exercise}

\begin{exercise}[3]
Explain detailed balance.
\end{exercise}

\begin{exercise}[3]
Explain why MCMC autocorrelation reduces effective sample size.
\end{exercise}

\begin{exercise}[3]
Explain the E-step and M-step of the EM algorithm.
\end{exercise}

\begin{exercise}[3]
Explain why EM likelihood values generally increase across iterations.
\end{exercise}

\begin{exercise}[3]
Design a simulation study to estimate the bias and MSE of two competing
estimators.
\end{exercise}

\begin{exercise}[3]
Explain how to estimate coverage probability through simulation.
\end{exercise}

\begin{exercise}[3]
Explain why random seeds matter for reproducibility.
\end{exercise}

\begin{exercise}[3]
Explain why parallel simulation requires independent random-number
streams.
\end{exercise}

\begin{exercise}[3]
Explain why sparse matrices can dramatically reduce memory usage.
\end{exercise}

\begin{exercise}[4]
Study convergence rates of gradient descent under strong convexity.
\end{exercise}

\begin{exercise}[4]
Derive quasi-Newton methods and the BFGS update.
\end{exercise}

\begin{exercise}[4]
Study trust-region optimization.
\end{exercise}

\begin{exercise}[4]
Investigate constrained optimization using Lagrange multipliers and
interior-point methods.
\end{exercise}

\begin{exercise}[4]
Study automatic differentiation using computational graphs.
\end{exercise}

\begin{exercise}[4]
Derive coordinate-descent updates for elastic-net regression.
\end{exercise}

\begin{exercise}[4]
Study proximal-gradient methods for lasso estimation.
\end{exercise}

\begin{exercise}[4]
Investigate stochastic-gradient convergence.
\end{exercise}

\begin{exercise}[4]
Study variance-reduced stochastic-gradient methods.
\end{exercise}

\begin{exercise}[4]
Derive acceptance--rejection efficiency.
\end{exercise}

\begin{exercise}[4]
Study importance-sampling central limit theorems.
\end{exercise}

\begin{exercise}[4]
Investigate adaptive importance sampling.
\end{exercise}

\begin{exercise}[4]
Study sequential Monte Carlo algorithms.
\end{exercise}

\begin{exercise}[4]
Derive particle-filter recursions.
\end{exercise}

\begin{exercise}[4]
Study theoretical validity of the nonparametric bootstrap.
\end{exercise}

\begin{exercise}[4]
Investigate bootstrap inconsistency for extreme-value statistics.
\end{exercise}

\begin{exercise}[4]
Study the block bootstrap for stationary time series.
\end{exercise}

\begin{exercise}[4]
Derive the BCa bootstrap interval.
\end{exercise}

\begin{exercise}[4]
Study permutation inference under restricted randomization schemes.
\end{exercise}

\begin{exercise}[4]
Derive Metropolis--Hastings invariance using detailed balance.
\end{exercise}

\begin{exercise}[4]
Study geometric ergodicity of Markov chains.
\end{exercise}

\begin{exercise}[4]
Investigate MCMC central limit theorems.
\end{exercise}

\begin{exercise}[4]
Study Hamiltonian Monte Carlo computationally.
\end{exercise}

\begin{exercise}[4]
Derive EM using Jensen's inequality.
\end{exercise}

\begin{exercise}[4]
Study acceleration methods for slow EM algorithms.
\end{exercise}

\begin{exercise}[4]
Investigate numerical linear algebra for very large least-squares
problems.
\end{exercise}

\begin{exercise}[4]
Study randomized low-rank matrix decompositions.
\end{exercise}

\begin{exercise}[4]
Investigate distributed statistical computation.
\end{exercise}

%%%%%%%%%%%%%%%%%%%%%%%%%%%%%%%%%%%%%%%%%%%%%%%%%%%%%%%%%%%%
\subsection{Conceptual Summary}
\label{subsec:computational-statistics-summary}
%%%%%%%%%%%%%%%%%%%%%%%%%%%%%%%%%%%%%%%%%%%%%%%%%%%%%%%%%%%%

Computational statistics provides algorithms for statistical problems
that cannot be solved conveniently in closed form.

Optimization solves problems such as

\[
\boxed{
\widehat{\theta}
=
\arg\max_\theta
\ell(\theta).
}
\]

Gradient descent uses

\[
\boxed{
\theta^{(k+1)}
=
\theta^{(k)}
-
\eta_k
\nabla Q
(
\theta^{(k)}
).
}
\]

Newton's method uses curvature:

\[
\boxed{
\theta^{(k+1)}
=
\theta^{(k)}
-
H^{-1}
\nabla Q.
}
\]

Monte Carlo approximates expectations by

\[
\boxed{
\widehat{\mu}
=
\frac1M
\sum_{m=1}^{M}
g(X_m).
}
\]

Its standard error decreases at rate

\[
\boxed{
M^{-1/2}.
}
\]

Bootstrap methods approximate sampling distributions by resampling.

Permutation methods approximate null distributions through valid
rearrangements.

MCMC constructs dependent draws from difficult target distributions.

EM alternates expected latent-data completion and maximization.

Stable statistical computation relies on:

\[
\boxed{
\text{numerical analysis}
+
\text{optimization}
+
\text{simulation}
+
\text{resampling}
+
\text{linear algebra}.
}
\]

%%%%%%%%%%%%%%%%%%%%%%%%%%%%%%%%%%%%%%%%%%%%%%%%%%%%%%%%%%%%
\subsection{Section Connections}
\label{subsec:computational-statistics-connections}
%%%%%%%%%%%%%%%%%%%%%%%%%%%%%%%%%%%%%%%%%%%%%%%%%%%%%%%%%%%%

\chapterconnection{
Computational Statistics completes Chapter~\ref{chap:statistics} by
providing the numerical machinery needed to implement the methods
developed throughout the chapter. Likelihood estimation requires
optimization, bootstrap and permutation methods provide computational
inference when analytic sampling distributions are unavailable, Bayesian
statistics relies heavily on Monte Carlo and MCMC, and statistical
learning requires scalable optimization and matrix algorithms. These
ideas connect directly with Chapter 10 on optimization and Chapter 11 on
numerical and computational mathematics, where the algorithms underlying
large-scale scientific computation are developed in greater depth.
}

%%%%%%%%%%%%%%%%%%%%%%%%%%%%%%%%%%%%%%%%%%%%%%%%%%%%%%%%%%%%
% SECTION SOURCE TRACKING
%%%%%%%%%%%%%%%%%%%%%%%%%%%%%%%%%%%%%%%%%%%%%%%%%%%%%%%%%%%%

%------------------------------------------------------------
% 8.19 Computational Statistics
%------------------------------------------------------------
%
% Source basis:
%   - Standard computational statistics.
%   - Standard numerical optimization for statistical estimation.
%   - Standard Monte Carlo, bootstrap, permutation, and MCMC methods.
%
% Used for:
%   - analytical versus computational inference;
%   - numerical error;
%   - floating-point arithmetic;
%   - numerical stability;
%   - log-likelihood computation;
%   - log-sum-exp stabilization;
%   - optimization;
%   - gradients and Hessians;
%   - gradient descent;
%   - Newton--Raphson;
%   - Fisher scoring;
%   - IRLS;
%   - convergence criteria;
%   - local and global optimization;
%   - constrained optimization;
%   - reparameterization;
%   - coordinate descent;
%   - soft thresholding;
%   - stochastic gradient descent;
%   - numerical differentiation;
%   - automatic differentiation preview;
%   - QR decomposition;
%   - singular value decomposition;
%   - condition numbers;
%   - Monte Carlo integration;
%   - Monte Carlo standard errors;
%   - pseudo-random number generation;
%   - inverse transform sampling;
%   - acceptance--rejection sampling;
%   - importance sampling;
%   - effective sample size;
%   - variance reduction;
%   - control variates;
%   - antithetic variables;
%   - nonparametric bootstrap;
%   - parametric bootstrap;
%   - bootstrap standard errors;
%   - bootstrap confidence intervals;
%   - BCa preview;
%   - dependent-data bootstrap;
%   - jackknife;
%   - permutation tests;
%   - randomization tests;
%   - Monte Carlo permutation tests;
%   - plus-one p-value correction;
%   - MCMC;
%   - Metropolis--Hastings;
%   - Gibbs sampling;
%   - MCMC diagnostics;
%   - sequential Monte Carlo preview;
%   - particle filtering preview;
%   - expectation--maximization;
%   - simulation studies;
%   - bias, MSE, coverage, Type I error, and power simulation;
%   - reproducibility;
%   - parallel computation;
%   - sparse matrices;
%   - computational complexity;
%   - low-rank approximation;
%   - high-dimensional computation;
%   - applications and exercises.
%
% Notes:
%   - This completes Chapter 8 — Statistics.
%   - Chapter 9 begins Differential Equations and Dynamical Systems.
%   - Chapter 10 develops optimization methods in greater depth.
%   - Chapter 11 develops numerical analysis and scientific computing in
%     greater depth.
%
%%%%%%%%%%%%%%%%%%%%%%%%%%%%%%%%%%%%%%%%%%%%%%%%%%%%%%%%%%%%